%% 第22回編集委員会（物理第2巻『熱学』）：組 熱W8
%% p2-15.tex  熱-15　並び方を数える：統計力学の入口（★3〜5・3段）
%% 計画書：tools/phys2-plan.md の B「熱-15」・A の N43・N57・C-1・C-6。数値はすべて checks/物理第2巻/熱-15/calc15.py で計算した（出力 calc15_out.txt）。
%% 記号（この章だけの分け方。記号カードに書いた）
%%   場合の数（ミクロな状態の数）は Ω（大文字のオメガ）。W は仕事の記号なので、この章では使わない（専門書の S＝k log W の W は Ω のこと、と専門書の囲みで読む）。
%%   この章では仕事を使わない。小文字の ω は角振動数（⑧、数学 6-6）。
%%   ボルツマン定数 k_B（いつも添字）、ばね定数 k_s。分配関数 Z。エネルギーの準位 ε_i、エネルギーの粒（準位の間かく）Δε。
%%   粒子の数 N、棚 i にいる数 n_i、確率 P_i（圧力 p と分けるため大文字）。エネルギーの粒の数 j（⑧の準位の番号も同じ j＝粒の数）。
%%   未定乗数は α・β（条件が2つなので2つ。数学 8-4 の μ は、この章では⑨の化学ポテンシャル μ に使うので、乗数には使わない）。β＝1/(k_B T)。
%%   内部エネルギー U、熱容量 C_th（N 個の熱容量）、モル比熱 C_V（1 mol あたり）、アインシュタイン温度 θ_E。高さ z、高さの目安 H＝k_BT/(mg)（熱-3 ⑥ と同じ式）。
%%   物質量はことばで「1 mol」と書き、n を物質量に使わない（n は「個数」：n! と n_i）。
%% 例題13（計画書の12に「2つの固体が粒を分け合う」を足した。章末問題7の形を例題で見せるため）。章末問題は基本7・チャレンジ5。
%% 図：はてな・やってみよう・図1〜15（並び方16通り、分布、けたの物差し、つぼ、粒と仕切り、2つの固体の山、スターリング、棚、ボルツマン因子、2準位、アインシュタイン、3つの分布）。
%% ほかの章との約束（計画書 C-4 の2）
%%   R ln2＝5.76 J/K（熱-6 ④ の自由膨張と同じ数）。等温の大気の高さの目安（熱-3 ⑥ と同じ形。窒素15 ℃で 8.7 km、空気で 8.4 km）。
%%   等分配 ½k_BT（熱-14 ④）、デュロン・プティの 3R（熱-14 ④）。マクスウェル分布（熱-14 ⑤）はボルツマン分布から、と⑥で受ける。
%%   数学 8-4 のもっと知りたい「物理では、決まった全体のエネルギーの中で、いちばん起こりやすい分け方をさがすときに使います（物理 熱-15）」を⑥で果たす。
%%   数学 3-17 のもっと知りたい「52枚の並び方 約8×10^67 通り」「気体の分子の並び方の数」を、はてな？と例題1で受ける。
%%   数学 7-10⑤ の「物理の本では x＝e^{−a} の形」を⑧で果たす。スターリングの近似は⑤で、5-6・5-10・5-11・5-14 を使って作る（数学の本にない道具）。
%% 出典（本文で使った事実・値）
%%   ・k_B＝1.380649×10^-23 J/K、N_A＝6.02214076×10^23 /mol、1 eV＝1.602176634×10^-19 J（2019年の SI の決めた値）。R＝N_A k_B＝8.314 J/(mol·K)。
%%   ・ボルツマン 1877年の論文：Wiener Berichte 76, 373–435 (1877)。英訳と解説 Entropy 17, 1971 (2015)（K. Sharp, F. Matschinsky）。
%%   ・S＝k log W の形と定数 k はプランク（1900〜1901年、熱ふく射の研究）。ウィーンの中央墓地のボルツマンの墓に S＝k. log W と刻まれている（複数の二次資料で確かめた）。
%%   ・エーレンフェストのつぼのモデル：P. & T. Ehrenfest, Physikalische Zeitschrift 8, 311–314 (1907)。全部が片方にもどるまでの平均の回数 2^N（M. Kac, 1947）。
%%   ・リフルシャッフル7回：D. Bayer, P. Diaconis, Ann. Appl. Probab. 2(2), 294–313 (1992)。
%%   ・ガラス転移：P. W. Anderson, "Through the Glass Lightly", Science 267, 1615 (1995)「固体の理論で最も深く、最も興味深い未解決の問題は、おそらくガラスとガラス転移の性質」。
%%   ・ランダウアーの原理：R. Landauer, IBM J. Res. Dev. 5, 183 (1961)。実験：A. Bérut ほか, Nature 483, 187 (2012)。シラード 1929年。
%%   ・アインシュタインの固体の比熱：A. Einstein, Ann. Phys. 22, 180 (1907)。デバイ 1912年。ショットキー（W. Schottky）1922年。
%%   ・ダイヤモンドのモル比熱 298 K で 6.1 J/(mol·K)：NIST WebBook の 6.115 を二次資料で確かめた〔要：WebBook の本文。webbook.nist.gov に届かなかった〕。θ_E≈1300 K はこの値に合うようにこの本で選んだ値（アインシュタインの論文の値ではない）。
%%   ・窒素の標準モルエントロピー 191.6 J/(mol·K)（25 ℃・1 bar）：化学便覧・NIST の表の値〔要：本文での照合〕。
%%   ・MRI（1.5 T、水素の原子核）の2つの準位の差：共鳴の周波数 42.58 MHz/T×1.5 T＝63.9 MHz から ΔE≈2.6×10^-7 eV。体温（310 K）で数の差は約 5×10^-6（RIT の教材 "The Basics of MRI" の問題の答え 4.94×10^-6 と一致）。
%%   ・アルミ電解コンデンサーの「10 ℃上がると寿命は約半分」：部品メーカーの技術資料の目安（アレニウスの式による）〔要：メーカーの資料の本文〕。
\makeatletter
% トランプのカード  \qxpcard{(中心)}{文字}{色}
\newcommand\qxpcard[3]{\begin{scope}[shift={#1}]
  \filldraw[fill=white, draw=bkInk!70, line width=0.5pt, rounded corners=0.7mm] (-3.2,-4.6) rectangle (3.2,4.6);
  \node[font=\sffamily\footnotesize, text=#3, inner sep=0pt] at (0,0.2) {#2};\end{scope}}
% 番号つきの玉  \qxpball{(中心)}{色}{番号}
\newcommand\qxpball[3]{\filldraw[fill=#2!35, draw=#2!80!black, line width=0.5pt] #1 circle (2.3);
  \node[font=\sffamily\footnotesize, inner sep=0pt] at #1 {#3};}
% 記号つきの分子（図の字は footnotesize）  \qxpm{(中心)}{色}{文字}
\newcommand\qxpm[3]{\filldraw[fill=#2!30, draw=#2!80!black, line width=0.45pt] #1 circle (1.8);
  \node[font=\sffamily\footnotesize, inner sep=0pt] at #1 {#3};}
% 小さな分子（番号なし）  \qxpmol{(中心)}{色}
\newcommand\qxpmol[2]{\filldraw[fill=#2!45, draw=#2!80!black, line width=0.4pt] #1 circle (1.15);}
% 2つに仕切った箱  \qxpbox{(左下)}{はば}{高さ}
\newcommand\qxpbox[3]{\begin{scope}[shift={#1}]
  \draw[bkInk, line width=0.6pt] (0,0) rectangle (#2,#3);
  \draw[bkGray, line width=0.4pt, dashed] (#2/2,0) -- (#2/2,#3);\end{scope}}
\makeatother
\tikzset{qxp lab/.style={font=\sffamily\footnotesize, inner sep=1pt},
  qxp arr/.style={-{Stealth[length=3.2pt, width=3pt]}, line width=0.8pt}}

\phchapter{15}{並び方を数える：\\統計力学の入口}{トランプは、切るほどなぜ順番がばらばらになる？}{\bklevelrange{3}{5}}{並び方を数える：統計力学の入口}
\bkmichishirube{3-17,5-14,5-17,7-10,8-4,熱-3,熱-6,熱-14}{熱-16,原-4,原-6,原-8,原-9}

\begin{goal}
\begin{itemize}[label={\color{bkTheme}●}, leftmargin=1.6\zw, itemsep=0pt, topsep=1pt]
\item \bklevel{3}\enspace 並び方の数を数えて、変化がなぜ一方の向きにだけ進むのかを説明できる。
\item \bklevel{3}\enspace $S=k_{\mathrm B}\ln\Omega$ で、並び方の数からエントロピーを計算できる。
\item \bklevel{5}\enspace ボルツマン分布と分配関数から、2準位の系や固体の比熱を導ける。
\end{itemize}
\end{goal}

\phkoushiki{%
\phk{$S=k_{\mathrm B}\ln\Omega$（ボルツマンの関係）}{高校物理（発展）}{ph:熱-15:boltzmann}
}

\begin{hatena}
\noindent\begin{minipage}[c]{0.53\linewidth}\raggedright
新しいトランプは、マークごとに A、2、3、…、K の順に並んでいます。何回か切ると、すぐにばらばらになります。ところが、ばらばらのトランプをいくら切っても、元の順にもどることは、まずありません。同じ「切る」なのに、なぜ「そろう→ばらばら」の向きにだけ進むのでしょう。
\end{minipage}\hfill
\begin{minipage}[c]{0.44\linewidth}\centering
\begin{tikzpicture}[x=1mm, y=1mm]
  \foreach \t [count=\i] in {A,2,3,4,5,6} {\qxpcard{(\i*7.6,28)}{\t}{bkInk}}
  \foreach \t/\c [count=\i] in {7/bkInk,K/bkRed,3/bkRed,J/bkInk,2/bkInk,9/bkRed} {\qxpcard{(\i*7.6,4)}{\t}{\c}}
  \draw[qxp arr, bkBlue] (19,22) -- node[left=1.2mm, qxp lab, text=bkBlue!80!black] {切る} (19,10);
  \draw[qxp arr, bkRed, dashed] (34,10) -- node[right=1.2mm, font=\sffamily\bfseries\large, text=bkRed] {?} (34,22);
  \node[qxp lab, anchor=south] at (26.6,33.5) {そろった順（$\spadesuit$）};
  \node[qxp lab, anchor=north] at (26.6,-1.5) {ばらばら};
\end{tikzpicture}
\end{minipage}
\end{hatena}

\noindent 熱-6 では、お湯と水をまぜると元にもどらないことを、熱÷温度を数えるエントロピー $S$ で表しました。熱-14 では、温度の正体が分子の運動のはげしさだとわかりました。この章では、2つをつなぎます。カギは「\textbf{並び方の数を数える}」ことです。分子は $10^{23}$ 個もあるので、並び方の数はとてつもなく大きくなります。その大きさが、熱の世界の「一方通行」を生んでいます。分子1つ1つの運動を追わずに、数と確率で熱をあつかう物理を\textbf{統計力学}といいます。\bkidx{とうけいりきがく}{統計力学}

\begin{yattemiyou}[さいころで、おはじきを動かそう]
\textbf{用意するもの}\par
さいころ1個、1〜6の番号を書いたおはじき6個（番号を書いた紙でもよい）、皿2枚\par\vspace{2pt}
\noindent\begin{minipage}[c]{0.55\linewidth}
\begin{enumerate}[label={\bkcircled[bkGreen]{\arabic*}}, leftmargin=2.4\zw, itemsep=1pt, topsep=0pt]
\item 6個のおはじきを、全部左の皿に入れる。
\item さいころを振り、出た目と同じ番号のおはじきを、反対の皿へ移す。
\item 左の皿にあるおはじきの数を記録する。②と③を30回繰り返す。
\end{enumerate}
\end{minipage}\hfill
\begin{minipage}[c]{0.42\linewidth}\centering
\begin{tikzpicture}[x=1mm, y=1mm]
  \draw[bkInk, line width=0.7pt] (0,4) arc (180:360:11 and 3.2);
  \draw[bkInk, line width=0.7pt] (0,4) -- (22,4);
  \draw[bkInk, line width=0.7pt] (28,4) arc (180:360:11 and 3.2);
  \draw[bkInk, line width=0.7pt] (28,4) -- (50,4);
  \foreach \x/\y/\n in {4/7/1,9/7/2,14/7/3,6.5/11.5/4,11.5/11.5/5,16.5/11.5/6} {\qxpball{(\x,\y)}{bkBlue}{\n}}
  \node[qxp lab] at (11,-2.5) {左の皿};
  \node[qxp lab] at (39,-2.5) {右の皿};
  \filldraw[fill=white, draw=bkInk, line width=0.6pt, rounded corners=0.8mm] (35,10) rectangle (42,17);
  \fill[bkInk] (36.8,15.2) circle (0.6) (38.5,13.5) circle (0.6) (40.2,11.8) circle (0.6);
  \draw[qxp arr, bkRed] (18,15) to[bend left=30] node[above, qxp lab, text=bkRed!80!black] {3番を移す} (30,15);
\end{tikzpicture}
\end{minipage}
\end{yattemiyou}

\noindent やってみようの結果の例です（図\ref{ph:熱-15:fig-saikoro}）。はじめは6個とも左にありましたが、左の数はすぐに減り、そのあとは3（ちょうど半分）のまわりを、大きく上がったり下がったりしました。6個とも左にもどったことも、4回め・18回め・22回めの3回ありました。6個のように少ないと、ばらつきが大きいのです。おはじきの数を増やすと、このふるまいがどう変わるかを、③で調べます。
\begin{bkfigure}
\begin{tikzpicture}
\begin{axis}[bkplot, width=118mm, height=40mm, xmin=0, xmax=31, ymin=0, ymax=6.8, xtick={5,10,15,20,25,30}, ytick={1,2,3,4,5,6}, xlabel={回}, ylabel={左の数}]
  \addplot[bkBlue, line width=1.1pt, const plot, mark=*, mark size=1.2pt, mark options={fill=bkBlue}] coordinates {(0,6) (1,5) (2,4) (3,5) (4,6) (5,5) (6,4) (7,3) (8,2) (9,3) (10,2) (11,1) (12,2) (13,3) (14,4) (15,5) (16,4) (17,5) (18,6) (19,5) (20,4) (21,5) (22,6) (23,5) (24,4) (25,3) (26,4) (27,5) (28,4) (29,3) (30,4)};
  \addplot[bkRed, line width=0.8pt, dashed, domain=0:31] {3};
  \node[font=\sffamily\footnotesize, text=bkRed!75!black, anchor=north east] at (axis cs:30.5,2.7) {ちょうど半分（3個）};
\end{axis}
\end{tikzpicture}
\bkfigcaption{さいころで動かしたおはじき6個の、左の皿の数（結果の例）。さいころの目は 2, 5, 5, 2, 3, 5, 4, 6, 5, 1, … でした。}\label{ph:熱-15:fig-saikoro}
\end{bkfigure}

\begin{omoidasou}
\bkrecall{3-17}{コイン $N$ 枚を投げて表が $k$ 枚になる並び方の数は、パスカルの三角形の $N$ 段めの数 ${}_N\mathrm{C}_k$。52枚のトランプの並べ方は約 $8\times10^{67}$ 通り。}
\bkrecall{5-17}{$n!=1\times2\times3\times\cdots\times n$（$n$ の階乗）。$0!=1$。}
\bkrecall{3-12}{対数はかけ算をたし算にする：$\log(ab)=\log a+\log b$。}
\bkrecall{5-14}{$\ln x$ は $e$ を底とする対数。$(\ln x)'=\tfrac{1}{x}$、$\int\tfrac{1}{x}\,\dd x=\ln x$。}
\bkrecall{熱-6}{可逆に熱 $Q$ を温度 $T$〔K〕で受けとると、エントロピーは $\Delta S=Q/T$ 増える。理想気体1 mol を等温で体積2倍にすると $\Delta S=R\ln2$。}
\bkrecall{熱-14}{温度は分子の運動のはげしさ。$\tfrac{1}{2}m\langle v^2\rangle=\tfrac{3}{2}k_{\mathrm B}T$、$k_{\mathrm B}=\qty{1.380649e-23}{J/K}$。}
\bkrecall{熱-3}{温度が一定の大気では、圧力は高さ $z$ とともに $p=p_0e^{-z/H}$ で下がる。}
\bkrecall{8-4}{条件つきの最大（ラグランジュの未定乗数法）：最大の点では、$\mathrm{grad}$ どうしが平行になる。}
\bkrecall{7-10}{等比級数：$|x|<1$ なら $1+x+x^2+\cdots=\dfrac{1}{1-x}$。}
\end{omoidasou}

%% ======================== 物理の道具 ========================
\Needspace{10\baselineskip}\bksection{doguu}{物理の道具}

\phstage{me}{並び方を数えて、変化の向きを見る}{3}

\bksubsection{1}{コインと分子：並び方の数を数える}
\label{ph:熱-15:sec1}
コイン4枚を投げます。1枚1枚の表・裏まで書くと、表表表表、表表表裏、…と $2^4=16$ 通りあります（3-17②）。ところが、表の枚数だけを見ると、0枚から4枚までの5通りしかありません。表が2枚になる並び方は6通りもありますが、4枚とも表になる並び方は1通りだけです。

次に、まん中に見えない線を引いた箱に、分子が4個（A・B・C・D）入っているとします。どの分子も、左の半分にも右の半分にもいられます。これは、コインの表・裏と同じです。16通りの並び方を、左にいる分子の数で分けると、図\ref{ph:熱-15:fig-16}のようになります。
\begin{bkfigure}
\begin{tikzpicture}[x=1mm, y=1mm]
  \foreach \x/\k/\w in {9/4/1,34/3/4,59/2/6,84/1/4,109/0/1} {\node[qxp lab, anchor=south, align=center] at (\x,10.5) {左に\k 個\\\w 通り};}
\qxpbox{(0,0)}{18}{9}
\qxpm{(2.185,6.6)}{bkBlue}{A}
\qxpm{(6.365,6.6)}{bkRed}{B}
\qxpm{(2.185,2.4)}{bkGreen}{C}
\qxpm{(6.365,2.4)}{bkOrange}{D}
\qxpbox{(25,0)}{18}{9}
\qxpm{(27.185,6.6)}{bkBlue}{A}
\qxpm{(31.365,6.6)}{bkRed}{B}
\qxpm{(27.185,2.4)}{bkGreen}{C}
\qxpm{(36.185,6.6)}{bkOrange}{D}
\qxpbox{(25,-10.4)}{18}{9}
\qxpm{(27.185,-3.8)}{bkBlue}{A}
\qxpm{(31.365,-3.8)}{bkRed}{B}
\qxpm{(27.185,-8)}{bkOrange}{D}
\qxpm{(36.185,-3.8)}{bkGreen}{C}
\qxpbox{(25,-20.8)}{18}{9}
\qxpm{(27.185,-14.2)}{bkBlue}{A}
\qxpm{(31.365,-14.2)}{bkGreen}{C}
\qxpm{(27.185,-18.4)}{bkOrange}{D}
\qxpm{(36.185,-14.2)}{bkRed}{B}
\qxpbox{(25,-31.2)}{18}{9}
\qxpm{(27.185,-24.6)}{bkRed}{B}
\qxpm{(31.365,-24.6)}{bkGreen}{C}
\qxpm{(27.185,-28.8)}{bkOrange}{D}
\qxpm{(36.185,-24.6)}{bkBlue}{A}
\qxpbox{(50,0)}{18}{9}
\qxpm{(52.185,6.6)}{bkBlue}{A}
\qxpm{(56.365,6.6)}{bkRed}{B}
\qxpm{(61.185,6.6)}{bkGreen}{C}
\qxpm{(65.365,6.6)}{bkOrange}{D}
\qxpbox{(50,-10.4)}{18}{9}
\qxpm{(52.185,-3.8)}{bkBlue}{A}
\qxpm{(56.365,-3.8)}{bkGreen}{C}
\qxpm{(61.185,-3.8)}{bkRed}{B}
\qxpm{(65.365,-3.8)}{bkOrange}{D}
\qxpbox{(50,-20.8)}{18}{9}
\qxpm{(52.185,-14.2)}{bkBlue}{A}
\qxpm{(56.365,-14.2)}{bkOrange}{D}
\qxpm{(61.185,-14.2)}{bkRed}{B}
\qxpm{(65.365,-14.2)}{bkGreen}{C}
\qxpbox{(50,-31.2)}{18}{9}
\qxpm{(52.185,-24.6)}{bkRed}{B}
\qxpm{(56.365,-24.6)}{bkGreen}{C}
\qxpm{(61.185,-24.6)}{bkBlue}{A}
\qxpm{(65.365,-24.6)}{bkOrange}{D}
\qxpbox{(50,-41.6)}{18}{9}
\qxpm{(52.185,-35)}{bkRed}{B}
\qxpm{(56.365,-35)}{bkOrange}{D}
\qxpm{(61.185,-35)}{bkBlue}{A}
\qxpm{(65.365,-35)}{bkGreen}{C}
\qxpbox{(50,-52)}{18}{9}
\qxpm{(52.185,-45.4)}{bkGreen}{C}
\qxpm{(56.365,-45.4)}{bkOrange}{D}
\qxpm{(61.185,-45.4)}{bkBlue}{A}
\qxpm{(65.365,-45.4)}{bkRed}{B}
\qxpbox{(75,0)}{18}{9}
\qxpm{(77.185,6.6)}{bkBlue}{A}
\qxpm{(86.185,6.6)}{bkRed}{B}
\qxpm{(90.365,6.6)}{bkGreen}{C}
\qxpm{(86.185,2.4)}{bkOrange}{D}
\qxpbox{(75,-10.4)}{18}{9}
\qxpm{(77.185,-3.8)}{bkRed}{B}
\qxpm{(86.185,-3.8)}{bkBlue}{A}
\qxpm{(90.365,-3.8)}{bkGreen}{C}
\qxpm{(86.185,-8)}{bkOrange}{D}
\qxpbox{(75,-20.8)}{18}{9}
\qxpm{(77.185,-14.2)}{bkGreen}{C}
\qxpm{(86.185,-14.2)}{bkBlue}{A}
\qxpm{(90.365,-14.2)}{bkRed}{B}
\qxpm{(86.185,-18.4)}{bkOrange}{D}
\qxpbox{(75,-31.2)}{18}{9}
\qxpm{(77.185,-24.6)}{bkOrange}{D}
\qxpm{(86.185,-24.6)}{bkBlue}{A}
\qxpm{(90.365,-24.6)}{bkRed}{B}
\qxpm{(86.185,-28.8)}{bkGreen}{C}
\qxpbox{(100,0)}{18}{9}
\qxpm{(111.185,6.6)}{bkBlue}{A}
\qxpm{(115.365,6.6)}{bkRed}{B}
\qxpm{(111.185,2.4)}{bkGreen}{C}
\qxpm{(115.365,2.4)}{bkOrange}{D}
\end{tikzpicture}
\bkfigcaption{分子4個の16通りの並び方。点線が箱のまん中。左に2個いる並び方がいちばん多い（6通り）。}\label{ph:熱-15:fig-16}
\end{bkfigure}

1つ1つの分子がどちらにいるかまで決めた並び方を\textbf{ミクロな状態}、「左に2個」のように全体の数だけで見た状態を\textbf{マクロな状態}といいます。\bkidx{みくろなじょうたい}{ミクロな状態}1つのマクロな状態に入るミクロな状態の数を、この本では $\Omega$（オメガ。ギリシャ文字の大文字）と書き、\textbf{状態の数}とよびます。図\ref{ph:熱-15:fig-16}では、$\Omega(4)=1$、$\Omega(3)=4$、$\Omega(2)=6$ です。専門書ではこの数を $W$ と書くことが多いのですが、この本では $W$ は仕事の記号なので、$\Omega$ を使います。

左に $k$ 個いる並び方の数は、$N$ 個の分子から、左にいる $k$ 個を選ぶ選び方の数（組合せ）なので、3-17②の ${}_N\mathrm{C}_k$ です。
\pheq{\Omega(k)={}_N\mathrm{C}_k=\frac{N!}{k!\,(N-k)!}}{通り ＝ 通り}
$N=4$、$k=2$ なら $\dfrac{4!}{2!\times2!}=\dfrac{24}{4}=6$ で、図と合います。

\phboxhead{どの並び方も、同じくらい起こる}
分子はとても速く飛び回り、壁や互いにぶつかって、向きを変え続けています（熱-14）。少したつと、どの分子がどちらにいるかは、まったくわからなくなります。そこで統計力学では、次のように考えます。
\begin{phmodel}[等確率の原理]
\phignore{分子1つ1つの細かい運動の道すじ。分子どうしが引き合う力（ここでは、分子は互いにじゃまをしない）。}
\phassume{全体のエネルギーなどの条件が同じなら、どのミクロな状態も、同じ確からしさで起こる（\textbf{等確率の原理}）。}
\phlimit{分子がなかなか動けず、いろいろな並び方をめぐれないとき（固まったガラスなど。⑨のあとの「まだわかっていないこと」）。}
\end{phmodel}
\bkidx{とうかくりつのげんり}{等確率の原理}この仮定のもとでは、左に $k$ 個いる確率 $P(k)$ は、状態の数を全部の並び方の数 $2^N$ でわったものです。
\pheq{P(k)=\frac{\Omega(k)}{2^N}}{1 ＝ 通り÷通り}
$N=4$ なら、左に2個の確率は $6\div16=0.375$、4個とも左の確率は $1\div16=0.0625$ です。4個しかないと、「全部が左に集まる」ことも、16回に1回は起こります。

\begin{phyomi}[$P(k)=\dfrac{\Omega(k)}{2^N}$]
「起こりやすいマクロな状態とは、ミクロな状態をたくさんもつ状態のこと」。どの並び方も同じくらい起こるので、並び方の数が多いマクロな状態ほど、よく見られます。
\end{phyomi}

\begin{phbaisuru}[10倍]{分子の数 $N$}
$N=4$・10・100 で、左にいる分子の割合 $k/N$ ごとの確率を並べました（図\ref{ph:熱-15:fig-bunpu}）。分子を増やすほど、山は「ちょうど半分」の所で細くとがります。左が40〜60\,\%に入る確率は、$N=10$ で0.66、$N=100$ で0.96です。$N=1000$ なら、45〜55\,\%に入る確率が0.9986になります。
\end{phbaisuru}
\begin{bkfigure}
\begin{tikzpicture}
\begin{axis}[bkplot, name=a, width=39mm, height=38mm, xmin=0, xmax=1.12, ymin=0, ymax=0.42, xtick={0,0.5,1}, ytick={0.1,0.2,0.3,0.4}, yticklabel style={/pgf/number format/fixed}, xlabel={$k/N$}, xlabel style={at={(axis description cs:0.5,0)}, anchor=north, yshift=-4.5mm}, ylabel={$P$}, title={\footnotesize $N=4$}, title style={at={(0.5,1.06)}}]
  \addplot[ycomb, bkBlue, line width=3pt] coordinates {(0,0.06250) (0.25,0.25000) (0.5,0.37500) (0.75,0.25000) (1,0.06250)};
\end{axis}
\begin{axis}[bkplot, at={(a.east)}, anchor=west, xshift=22mm, name=b, width=39mm, height=38mm, xmin=0, xmax=1.12, ymin=0, ymax=0.28, xtick={0,0.5,1}, ytick={0.1,0.2}, yticklabel style={/pgf/number format/fixed}, xlabel={$k/N$}, xlabel style={at={(axis description cs:0.5,0)}, anchor=north, yshift=-4.5mm}, ylabel={$P$}, title={\footnotesize $N=10$}, title style={at={(0.5,1.06)}}]
  \addplot[ycomb, bkBlue, line width=2pt] coordinates {(0,0.00098) (0.1,0.00977) (0.2,0.04395) (0.3,0.11719) (0.4,0.20508) (0.5,0.24609) (0.6,0.20508) (0.7,0.11719) (0.8,0.04395) (0.9,0.00977) (1,0.00098)};
\end{axis}
\begin{axis}[bkplot, at={(b.east)}, anchor=west, xshift=22mm, width=39mm, height=38mm, xmin=0, xmax=1.12, ymin=0, ymax=0.09, xtick={0,0.5,1}, ytick={0.04,0.08}, yticklabel style={/pgf/number format/fixed}, xlabel={$k/N$}, xlabel style={at={(axis description cs:0.5,0)}, anchor=north, yshift=-4.5mm}, ylabel={$P$}, title={\footnotesize $N=100$}, title style={at={(0.5,1.06)}}]
  \addplot[ycomb, bkBlue, line width=0.5pt] coordinates {(0,0.00000) (0.01,0.00000) (0.02,0.00000) (0.03,0.00000) (0.04,0.00000) (0.05,0.00000) (0.06,0.00000) (0.07,0.00000) (0.08,0.00000) (0.09,0.00000) (0.1,0.00000) (0.11,0.00000) (0.12,0.00000) (0.13,0.00000) (0.14,0.00000) (0.15,0.00000) (0.16,0.00000) (0.17,0.00000) (0.18,0.00000) (0.19,0.00000) (0.2,0.00000) (0.21,0.00000) (0.22,0.00000) (0.23,0.00000) (0.24,0.00000) (0.25,0.00000) (0.26,0.00000) (0.27,0.00000) (0.28,0.00000) (0.29,0.00001) (0.3,0.00002) (0.31,0.00005) (0.32,0.00011) (0.33,0.00023) (0.34,0.00046) (0.35,0.00086) (0.36,0.00156) (0.37,0.00270) (0.38,0.00447) (0.39,0.00711) (0.4,0.01084) (0.41,0.01587) (0.42,0.02229) (0.43,0.03007) (0.44,0.03895) (0.45,0.04847) (0.46,0.05796) (0.47,0.06659) (0.48,0.07353) (0.49,0.07803) (0.5,0.07959) (0.51,0.07803) (0.52,0.07353) (0.53,0.06659) (0.54,0.05796) (0.55,0.04847) (0.56,0.03895) (0.57,0.03007) (0.58,0.02229) (0.59,0.01587) (0.6,0.01084) (0.61,0.00711) (0.62,0.00447) (0.63,0.00270) (0.64,0.00156) (0.65,0.00086) (0.66,0.00046) (0.67,0.00023) (0.68,0.00011) (0.69,0.00005) (0.7,0.00002) (0.71,0.00001) (0.72,0.00000) (0.73,0.00000) (0.74,0.00000) (0.75,0.00000) (0.76,0.00000) (0.77,0.00000) (0.78,0.00000) (0.79,0.00000) (0.8,0.00000) (0.81,0.00000) (0.82,0.00000) (0.83,0.00000) (0.84,0.00000) (0.85,0.00000) (0.86,0.00000) (0.87,0.00000) (0.88,0.00000) (0.89,0.00000) (0.9,0.00000) (0.91,0.00000) (0.92,0.00000) (0.93,0.00000) (0.94,0.00000) (0.95,0.00000) (0.96,0.00000) (0.97,0.00000) (0.98,0.00000) (0.99,0.00000) (1,0.00000)};
\end{axis}
\end{tikzpicture}
\bkfigcaption{左にいる分子の割合 $k/N$ と、その確率 $P$。たての目もりは3つのグラフでちがう。$N$ を増やすほど、半分のまわりに集まる。}\label{ph:熱-15:fig-bunpu}
\end{bkfigure}

\bksubsection{2}{並び方の数の爆発と、はてな？の答え}
\label{ph:熱-15:sec2}
分子の数を増やしていくと、並び方の数はどうなるでしょう。全部の並び方の数 $2^N$、ちょうど半分ずつになる並び方の数 ${}_N\mathrm{C}_{N/2}$、全部が左に集まる確率 $2^{-N}$ を表にしました。
\begin{bkcenter}
\begin{bktable}{colspec={Q[c,m,wd=20mm] Q[c,m,wd=28mm] Q[c,m,wd=30mm] Q[c,m,wd=32mm]}, rowsep=1.5pt, row{1}={bg=bkTheme!10}}
分子の数 $N$ & 全部の並び方 $2^N$ & 半分ずつの並び方 & 全部が左の確率 \\
4 & 16 & 6 & 0.0625 \\
10 & \num{1024} & 252 & \num{9.8e-4} \\
100 & \num{1.3e30} & \num{1.0e29} & \num{7.9e-31} \\
\num{1000} & \num{1.1e301} & \num{2.7e299} & \num{9.3e-302} \\
\end{bktable}
\bktabcaption{分子が増えると、並び方の数は爆発的に大きくなり、「全部が左」の確率は爆発的に小さくなる。}\label{ph:熱-15:tab-bakuhatsu}
\end{bkcenter}
分子が1つ増えるたびに、全部の並び方は2倍になります。このように、数が増えるたびに同じ割合でかけ算されて増えることを\textbf{指数関数的に増える}といいます（3-11）。1 mol（約 $6\times10^{23}$ 個）の気体なら、全部が左に集まる確率は、10を $1.8\times10^{23}$ 回かけた数で1をわったもの、$10^{-1.8\times10^{23}}$ です。小数点のあとに続く0を1つ\qty{2}{mm}で書いても、0の列は約 $4\times10^{20}$ m、光でも約4万年かかる長さになります。

\begin{phbaisuru}{分子の数 $N$}
全部の並び方の数 $2^N$ は、$2^{2N}=(2^N)^2$ なので\textbf{2乗}になります。100個で $1.3\times10^{30}$ なら、200個では約 $1.6\times10^{60}$ です。一方、半分からのずれの「ふつうの大きさ」（③で、およそ $\sqrt N/2$ 個だと確かめます）は $\sqrt2$ 倍にしかなりません。だから、全体の数に対するずれの割合は $1/\sqrt2$ 倍に小さくなります。
\end{phbaisuru}

\phboxhead{はてな？の答え}
\label{ph:熱-15:hatena}\bkhatenaref に答えましょう。52枚のトランプの並べ方は $52!$ 通りで、計算すると約 $8.07\times10^{67}$ 通りです（例題1。3-17のもっと知りたい）。このうち「きれいにそろった並び」は、ほんのわずかです。たとえば「4つのマークが、それぞれ A から K の順に、マークごとにかたまって並ぶ」並び方は、マークの順番の分だけで $4!=24$ 通りしかありません。その割合は
\[\frac{24}{8.07\times10^{67}}\approx3\times10^{-67}\]
です。52!がどれほど大きいかを、ほかの大きな数と、けた（10の何乗か）でくらべました（図\ref{ph:熱-15:fig-keta}）。
\begin{bkfigure}
\begin{tikzpicture}[x=1.7mm, y=1mm]
  \draw[bkInk, line width=0.8pt, -{Stealth[length=4pt]}] (0,0) -- (72,0);
  \foreach \k in {0,10,...,70} {\draw[bkInk, line width=0.6pt] (\k,-1.2) -- (\k,1.2); \node[qxp lab, anchor=north] at (\k,-1.8) {$10^{\k}$};}
  \foreach \x/\a/\t/\c in {17.6/south east/{宇宙の年齢の秒数\\$4\times10^{17}$}/bkGreen!60!black, 30.1/south west/{コイン100枚の並び方\\$1.3\times10^{30}$}/bkOrange!80!black, 67.9/south/{52枚の並べ方\\$8\times10^{67}$}/bkRed!80!black} {
    \fill[\c] (\x,0) circle[radius=1mm]; \node[qxp lab, align=center, text=\c, anchor=\a] at (\x,2.2) {\t};}
  \fill[bkBlue!80!black] (23.8,0) circle[radius=1mm];
  \node[qxp lab, align=center, text=bkBlue!80!black, anchor=north] at (23.8,-6.5) {1 mol の分子\\$6\times10^{23}$};
\end{tikzpicture}
\bkfigcaption{大きな数を、けたでくらべる。目もり1つで10の10乗倍。52枚のトランプの並べ方は、1 mol の分子の数より44けたも大きい。}\label{ph:熱-15:fig-keta}
\end{bkfigure}
よく切ったトランプは、52!通りのどの並び方にも同じくらいなりやすくなります（例題12）。すると、切ったあとに「ばらばら」になる確率はほとんど1、「そろう」確率はほとんど0です。

つまり、切るという動作が、ばらばらにする向きを好んでいるのではありません。\textbf{ばらばらな並び方が、そろった並び方より、けたちがいに多い}だけなのです。そろった並びにもどることは、禁じられてはいません。ただ、1秒に1回ずつ切り続けて、宇宙が生まれてから今まで（約 $4\times10^{17}$ 秒）続けても、全部の並び方の $10^{-50}$ ほどしかためせないので、まず起こらないのです。

\begin{phmijika}[インクが広がる・お湯と水がまざる]
水にインクを1滴落とすと、インクの分子は広がっていき、1か所に集まりもどることはありません。熱-6 で見た「もどらない変化」は、どれも同じです。分子1つ1つの運動は、時間を逆にしても成り立つ動きです（力-2）。それでも、広がった並び方のほうが、集まった並び方より、けたちがいに多いので、分子の数が多いと、もどる向きは事実上起こりません。これが、熱の世界の「一方通行」の正体です。
\end{phmijika}

\Needspace{16\baselineskip}\phstage{te}{表計算とプログラムで動かす}{3}

\bksubsection{3}{エーレンフェストのつぼ：表計算で動かす}
\label{ph:熱-15:sec3}
やってみようのゲームを、玉の数を増やして、表計算で行います。このゲームは、エーレンフェスト夫妻（パウルとタチアナ）が1907年に、エントロピーが増えるわけを説明するために考えたもので、\textbf{エーレンフェストのつぼ}とよばれます。\bkidx{えーれんふぇすとのつぼ}{エーレンフェストのつぼ}

\noindent\begin{minipage}[c]{0.5\linewidth}\raggedright
\begin{enumerate}[label={\bkcircled[bkBlue]{\arabic*}}, leftmargin=2.4\zw, itemsep=1pt, topsep=0pt]
\item 番号のついた玉 $N$ 個を、全部左のつぼに入れる。
\item 毎回、$N$ 個の中から1個を同じ確からしさで選び、反対のつぼへ移す。
\item 左に $L$ 個あるとき、選んだ玉が左にある確率は $L/N$。そのときは $L$ が1減り、そうでなければ1増える。
\end{enumerate}
\end{minipage}\hfill
\begin{minipage}[c]{0.46\linewidth}\centering
\begin{tikzpicture}[x=1mm, y=1mm]
  \draw[bkInk, line width=0.7pt, rounded corners=3mm] (0,22) -- (2,2) -- (20,2) -- (22,22);
  \draw[bkInk, line width=0.7pt, rounded corners=3mm] (32,22) -- (34,2) -- (52,2) -- (54,22);
  \foreach \x/\y/\n in {6/6/1,11/6/2,16/6/3,8.5/10.8/4,13.5/10.8/5,11/15.6/6} {\qxpball{(\x,\y)}{bkBlue}{\n}}
  \foreach \x/\y/\n in {40/6/7,45/6/8} {\qxpball{(\x,\y)}{bkBlue}{\n}}
  \draw[qxp arr, bkRed] (13.5,19) to[bend left=35] node[above, qxp lab, text=bkRed!80!black] {選んだ1個を移す} (41,19);
  \node[qxp lab] at (11,-2) {左：$L$ 個};
  \node[qxp lab] at (43,-2) {右：$N-L$ 個};
\end{tikzpicture}
\end{minipage}

\par\vspace{4pt}
\noindent\begin{minipage}[c]{0.5\linewidth}\centering
\bksheet{3}{16}{%
1 & 回 & RAND() & 左の数 \\
2 & 0 &  & 10 \\
3 & 1 & 0.11 & 9 \\
4 & 2 & 0.70 & 8 \\
5 & 3 & 0.65 & 7 \\
6 & 4 & 0.94 & 8 \\
7 & 5 & 0.27 & 7 \\
8 & 6 & 0.26 & 6 \\
9 & 7 & 0.73 & 7 \\
10 & 8 & 0.66 & 6 \\
11 & 9 & 0.30 & 5 \\
12 & 10 & 0.68 & 6 \\}
\bktabcaption{$N=10$ のシートの、はじめの10回（1回の例）。}\label{ph:熱-15:tab-sheet}
\end{minipage}\hfill
\begin{minipage}[c]{0.46\linewidth}\raggedright
表計算の RAND() は、0から1までの数を同じ確からしさで出す関数です。「RAND() の値が $L/N$ より小さい」ことを「左の玉を選んだ」ことにします。\par\vspace{3pt}
B3 に \texttt{=RAND()}、C3 に\par
\texttt{=IF(B3<C2/10,C2-1,C2+1)}\par
と入れて、下へコピーします。1回め（3行め）は、$0.11<10/10$ なので左の玉を選び、$L$ は9になりました。4回めは $0.94\geqq7/10$ なので、右の玉を選び、$L$ は8にもどりました。RAND の値は、計算し直すたびに変わります。
\end{minipage}

\par\vspace{4pt}
同じシートで、$N=10$・100・\num{1000} のときを\num{2000}回まで動かし、左にある割合 $L/N$ をグラフにしました（図\ref{ph:熱-15:fig-tsubo}）。
\begin{bkfigure}
\begin{tikzpicture}
\begin{axis}[bkplot, scale only axis, width=96mm, height=15mm, xmin=0, xmax=2080, ymin=0, ymax=1.15, xtick={500,1000,1500,2000}, ytick={0,0.5,1}, yticklabel style={/pgf/number format/fixed}, xticklabels={}, ylabel={$L/N$}]
  \addplot[bkGray!80!black, line width=0.5pt] coordinates {(0,1.000) (2,0.800) (4,0.600) (6,0.400) (8,0.400) (10,0.400) (12,0.200) (14,0.400) (16,0.200) (18,0.400) (20,0.400) (22,0.400) (24,0.400) (26,0.200) (28,0.400) (30,0.400) (32,0.400) (34,0.400) (36,0.200) (38,0.400) (40,0.400) (42,0.400) (44,0.600) (46,0.600) (48,0.600) (50,0.600) (52,0.600) (54,0.600) (56,0.400) (58,0.400) (60,0.600) (62,0.600) (64,0.600) (66,0.400) (68,0.600) (70,0.600) (72,0.600) (74,0.800) (76,0.800) (78,0.600) (80,0.400) (82,0.200) (84,0.200) (86,0.200) (88,0.400) (90,0.400) (92,0.600) (94,0.800) (96,0.600) (98,0.600) (100,0.600) (102,0.400) (104,0.400) (106,0.400) (108,0.400) (110,0.400) (112,0.600) (114,0.600) (116,0.600) (118,0.800) (120,0.800) (122,0.600) (124,0.600) (126,0.400) (128,0.200) (130,0.200) (132,0.200) (134,0.200) (136,0.200) (138,0.200) (140,0.400) (142,0.200) (144,0.400) (146,0.400) (148,0.200) (150,0.200) (152,0.200) (154,0.200) (156,0.200) (158,0.200) (160,0.400) (162,0.600) (164,0.400) (166,0.400) (168,0.600) (170,0.400) (172,0.600) (174,0.800) (176,0.800) (178,0.600) (180,0.400) (182,0.200) (184,0.400) (186,0.400) (188,0.600) (190,0.600) (192,0.400) (194,0.200) (196,0.400) (198,0.400) (200,0.600) (202,0.600) (204,0.800) (206,0.600) (208,0.600) (210,0.600) (212,0.600) (214,0.600) (216,0.600) (218,0.400) (220,0.600) (222,0.600) (224,0.600) (226,0.600) (228,0.400) (230,0.600) (232,0.600) (234,0.600) (236,0.600) (238,0.400) (240,0.400) (242,0.400) (244,0.400) (246,0.400) (248,0.600) (250,0.600) (252,0.600) (254,0.400) (256,0.400) (258,0.400) (260,0.400) (262,0.600) (264,0.600) (266,0.600) (268,0.600) (270,0.400) (272,0.600) (274,0.600) (276,0.600) (278,0.600) (280,0.600) (282,0.600) (284,0.600) (286,0.800) (288,0.800) (290,0.800) (292,0.800) (294,0.600) (296,0.400) (298,0.400) (300,0.600) (302,0.600) (304,0.600) (306,0.800) (308,0.800) (310,0.600) (312,0.800) (314,0.600) (316,0.400) (318,0.400) (320,0.400) (322,0.400) (324,0.400) (326,0.600) (328,0.600) (330,0.800) (332,0.600) (334,0.600) (336,0.400) (338,0.600) (340,0.600) (342,0.600) (344,0.600) (346,0.400) (348,0.400) (350,0.400) (352,0.600) (354,0.600) (356,0.600) (358,0.400) (360,0.600) (362,0.400) (364,0.400) (366,0.200) (368,0.200) (370,0.400) (372,0.400) (374,0.400) (376,0.200) (378,0.200) (380,0.400) (382,0.400) (384,0.400) (386,0.400) (388,0.600) (390,0.600) (392,0.800) (394,0.800) (396,0.600) (398,0.400) (400,0.200) (402,0.200) (404,0.200) (406,0.200) (408,0.400) (410,0.200) (412,0.400) (414,0.400) (416,0.400) (418,0.600) (420,0.600) (422,0.400) (424,0.400) (426,0.600) (428,0.600) (430,0.400) (432,0.400) (434,0.200) (436,0.200) (438,0.400) (440,0.600) (442,0.800) (444,0.600) (446,0.600) (448,0.600) (450,0.600) (452,0.600) (454,0.800) (456,0.600) (458,0.600) (460,0.800) (462,0.800) (464,0.600) (466,0.400) (468,0.400) (470,0.400) (472,0.600) (474,0.600) (476,0.400) (478,0.600) (480,0.600) (482,0.600) (484,0.600) (486,0.400) (488,0.400) (490,0.600) (492,0.400) (494,0.600) (496,0.600) (498,0.600) (500,0.400) (502,0.400) (504,0.400) (506,0.200) (508,0.400) (510,0.400) (512,0.400) (514,0.400) (516,0.400) (518,0.600) (520,0.400) (522,0.600) (524,0.400) (526,0.400) (528,0.400) (530,0.400) (532,0.400) (534,0.600) (536,0.800) (538,0.800) (540,0.600) (542,0.400) (544,0.400) (546,0.400) (548,0.400) (550,0.400) (552,0.400) (554,0.600) (556,0.600) (558,0.400) (560,0.600) (562,0.400) (564,0.200) (566,0.200) (568,0.400) (570,0.200) (572,0.200) (574,0.400) (576,0.600) (578,0.800) (580,0.800) (582,0.600) (584,0.600) (586,0.600) (588,0.400) (590,0.400) (592,0.400) (594,0.400) (596,0.400) (598,0.400) (600,0.400) (602,0.200) (604,0.400) (606,0.400) (608,0.600) (610,0.400) (612,0.200) (614,0.200) (616,0.200) (618,0.400) (620,0.400) (622,0.600) (624,0.400) (626,0.200) (628,0.200) (630,0.400) (632,0.400) (634,0.200) (636,0.400) (638,0.600) (640,0.600) (642,0.600) (644,0.600) (646,0.400) (648,0.400) (650,0.600) (652,0.600) (654,0.400) (656,0.600) (658,0.800) (660,0.800) (662,0.600) (664,0.600) (666,0.600) (668,0.600) (670,0.400) (672,0.600) (674,0.400) (676,0.600) (678,0.400) (680,0.400) (682,0.400) (684,0.600) (686,0.600) (688,0.400) (690,0.600) (692,0.400) (694,0.600) (696,0.400) (698,0.600) (700,0.400) (702,0.400) (704,0.400) (706,0.600) (708,0.400) (710,0.400) (712,0.400) (714,0.400) (716,0.400) (718,0.400) (720,0.400) (722,0.600) (724,0.400) (726,0.400) (728,0.200) (730,0.200) (732,0.400) (734,0.600) (736,0.600) (738,0.600) (740,0.600) (742,0.600) (744,0.400) (746,0.200) (748,0.400) (750,0.400) (752,0.400) (754,0.400) (756,0.400) (758,0.400) (760,0.400) (762,0.400) (764,0.400) (766,0.600) (768,0.600) (770,0.400) (772,0.400) (774,0.600) (776,0.400) (778,0.600) (780,0.600) (782,0.400) (784,0.400) (786,0.600) (788,0.600) (790,0.400) (792,0.600) (794,0.400) (796,0.400) (798,0.400) (800,0.600) (802,0.800) (804,0.600) (806,0.400) (808,0.400) (810,0.400) (812,0.200) (814,0.200) (816,0.400) (818,0.600) (820,0.400) (822,0.400) (824,0.400) (826,0.400) (828,0.600) (830,0.800) (832,0.600) (834,0.600) (836,0.400) (838,0.400) (840,0.200) (842,0.200) (844,0.400) (846,0.400) (848,0.400) (850,0.600) (852,0.600) (854,0.600) (856,0.800) (858,0.600) (860,0.400) (862,0.600) (864,0.600) (866,0.600) (868,0.600) (870,0.600) (872,0.600) (874,0.400) (876,0.200) (878,0.400) (880,0.400) (882,0.400) (884,0.400) (886,0.200) (888,0.400) (890,0.200) (892,0.400) (894,0.400) (896,0.400) (898,0.400) (900,0.400) (902,0.400) (904,0.400) (906,0.600) (908,0.400) (910,0.400) (912,0.400) (914,0.600) (916,0.600) (918,0.400) (920,0.400) (922,0.600) (924,0.600) (926,0.600) (928,0.400) (930,0.600) (932,0.400) (934,0.600) (936,0.400) (938,0.600) (940,0.400) (942,0.400) (944,0.600) (946,0.600) (948,0.600) (950,0.600) (952,0.600) (954,0.400) (956,0.400) (958,0.200) (960,0.400) (962,0.400) (964,0.400) (966,0.400) (968,0.600) (970,0.600) (972,0.600) (974,0.600) (976,0.400) (978,0.400) (980,0.400) (982,0.600) (984,0.400) (986,0.400) (988,0.400) (990,0.400) (992,0.400) (994,0.400) (996,0.600) (998,0.400) (1000,0.400) (1002,0.400) (1004,0.600) (1006,0.800) (1008,0.800) (1010,0.800) (1012,0.800) (1014,0.800) (1016,0.600) (1018,0.600) (1020,0.600) (1022,0.600) (1024,0.600) (1026,0.400) (1028,0.400) (1030,0.200) (1032,0.200) (1034,0.400) (1036,0.400) (1038,0.400) (1040,0.400) (1042,0.600) (1044,0.600) (1046,0.600) (1048,0.600) (1050,0.600) (1052,0.600) (1054,0.600) (1056,0.400) (1058,0.400) (1060,0.400) (1062,0.600) (1064,0.600) (1066,0.400) (1068,0.200) (1070,0.400) (1072,0.600) (1074,0.400) (1076,0.400) (1078,0.400) (1080,0.400) (1082,0.400) (1084,0.400) (1086,0.400) (1088,0.400) (1090,0.200) (1092,0.200) (1094,0.400) (1096,0.400) (1098,0.400) (1100,0.400) (1102,0.400) (1104,0.600) (1106,0.600) (1108,0.600) (1110,0.400) (1112,0.400) (1114,0.400) (1116,0.400) (1118,0.400) (1120,0.600) (1122,0.600) (1124,0.400) (1126,0.400) (1128,0.600) (1130,0.600) (1132,0.600) (1134,0.600) (1136,0.400) (1138,0.200) (1140,0.400) (1142,0.600) (1144,0.600) (1146,0.400) (1148,0.600) (1150,0.600) (1152,0.600) (1154,0.600) (1156,0.400) (1158,0.600) (1160,0.600) (1162,0.600) (1164,0.600) (1166,0.600) (1168,0.800) (1170,0.800) (1172,0.800) (1174,0.800) (1176,0.600) (1178,0.600) (1180,0.600) (1182,0.400) (1184,0.400) (1186,0.200) (1188,0.000) (1190,0.200) (1192,0.400) (1194,0.400) (1196,0.400) (1198,0.600) (1200,0.600) (1202,0.800) (1204,0.600) (1206,0.400) (1208,0.600) (1210,0.800) (1212,0.600) (1214,0.400) (1216,0.400) (1218,0.400) (1220,0.200) (1222,0.400) (1224,0.400) (1226,0.600) (1228,0.600) (1230,0.400) (1232,0.600) (1234,0.600) (1236,0.400) (1238,0.400) (1240,0.400) (1242,0.200) (1244,0.400) (1246,0.400) (1248,0.600) (1250,0.400) (1252,0.400) (1254,0.400) (1256,0.200) (1258,0.400) (1260,0.400) (1262,0.400) (1264,0.600) (1266,0.600) (1268,0.400) (1270,0.400) (1272,0.400) (1274,0.200) (1276,0.400) (1278,0.400) (1280,0.400) (1282,0.600) (1284,0.600) (1286,0.400) (1288,0.400) (1290,0.600) (1292,0.400) (1294,0.600) (1296,0.400) (1298,0.400) (1300,0.400) (1302,0.600) (1304,0.400) (1306,0.400) (1308,0.400) (1310,0.200) (1312,0.400) (1314,0.400) (1316,0.400) (1318,0.600) (1320,0.600) (1322,0.600) (1324,0.400) (1326,0.200) (1328,0.400) (1330,0.600) (1332,0.600) (1334,0.600) (1336,0.600) (1338,0.600) (1340,0.600) (1342,0.400) (1344,0.600) (1346,0.800) (1348,0.800) (1350,0.600) (1352,0.600) (1354,0.400) (1356,0.600) (1358,0.600) (1360,0.400) (1362,0.400) (1364,0.600) (1366,0.600) (1368,0.600) (1370,0.600) (1372,0.600) (1374,0.600) (1376,0.600) (1378,0.600) (1380,0.600) (1382,0.600) (1384,0.600) (1386,0.600) (1388,0.600) (1390,0.600) (1392,0.600) (1394,0.400) (1396,0.400) (1398,0.400) (1400,0.600) (1402,0.400) (1404,0.400) (1406,0.600) (1408,0.600) (1410,0.400) (1412,0.400) (1414,0.600) (1416,0.600) (1418,0.600) (1420,0.600) (1422,0.800) (1424,0.600) (1426,0.400) (1428,0.400) (1430,0.400) (1432,0.200) (1434,0.200) (1436,0.200) (1438,0.200) (1440,0.400) (1442,0.600) (1444,0.600) (1446,0.800) (1448,0.800) (1450,0.800) (1452,0.800) (1454,0.600) (1456,0.600) (1458,0.600) (1460,0.400) (1462,0.400) (1464,0.200) (1466,0.400) (1468,0.400) (1470,0.600) (1472,0.800) (1474,0.800) (1476,0.600) (1478,0.600) (1480,0.600) (1482,0.400) (1484,0.200) (1486,0.200) (1488,0.200) (1490,0.400) (1492,0.600) (1494,0.600) (1496,0.600) (1498,0.800) (1500,0.600) (1502,0.600) (1504,0.600) (1506,0.600) (1508,0.800) (1510,0.800) (1512,0.600) (1514,0.600) (1516,0.600) (1518,0.600) (1520,0.400) (1522,0.600) (1524,0.400) (1526,0.600) (1528,0.600) (1530,0.600) (1532,0.400) (1534,0.600) (1536,0.800) (1538,0.600) (1540,0.600) (1542,0.800) (1544,0.800) (1546,0.800) (1548,0.600) (1550,0.400) (1552,0.600) (1554,0.800) (1556,0.800) (1558,0.600) (1560,0.600) (1562,0.600) (1564,0.600) (1566,0.600) (1568,0.400) (1570,0.600) (1572,0.600) (1574,0.800) (1576,0.600) (1578,0.400) (1580,0.400) (1582,0.600) (1584,0.600) (1586,0.400) (1588,0.600) (1590,0.600) (1592,0.400) (1594,0.600) (1596,0.600) (1598,0.400) (1600,0.400) (1602,0.200) (1604,0.200) (1606,0.200) (1608,0.200) (1610,0.000) (1612,0.200) (1614,0.200) (1616,0.200) (1618,0.400) (1620,0.400) (1622,0.600) (1624,0.600) (1626,0.600) (1628,0.600) (1630,0.600) (1632,0.600) (1634,0.600) (1636,0.600) (1638,0.400) (1640,0.400) (1642,0.600) (1644,0.600) (1646,0.600) (1648,0.600) (1650,0.600) (1652,0.800) (1654,0.600) (1656,0.600) (1658,0.800) (1660,0.600) (1662,0.800) (1664,0.600) (1666,0.600) (1668,0.600) (1670,0.800) (1672,0.800) (1674,0.600) (1676,0.600) (1678,0.800) (1680,0.800) (1682,0.800) (1684,0.600) (1686,0.800) (1688,0.800) (1690,0.600) (1692,0.600) (1694,0.600) (1696,0.600) (1698,0.600) (1700,0.400) (1702,0.400) (1704,0.400) (1706,0.400) (1708,0.400) (1710,0.600) (1712,0.600) (1714,0.600) (1716,0.400) (1718,0.400) (1720,0.400) (1722,0.600) (1724,0.400) (1726,0.400) (1728,0.400) (1730,0.400) (1732,0.400) (1734,0.400) (1736,0.400) (1738,0.400) (1740,0.400) (1742,0.400) (1744,0.600) (1746,0.800) (1748,0.600) (1750,0.400) (1752,0.200) (1754,0.400) (1756,0.600) (1758,0.600) (1760,0.400) (1762,0.400) (1764,0.400) (1766,0.600) (1768,0.600) (1770,0.400) (1772,0.600) (1774,0.400) (1776,0.400) (1778,0.600) (1780,0.600) (1782,0.800) (1784,0.800) (1786,0.600) (1788,0.800) (1790,0.600) (1792,0.600) (1794,0.400) (1796,0.200) (1798,0.200) (1800,0.200) (1802,0.200) (1804,0.200) (1806,0.200) (1808,0.400) (1810,0.600) (1812,0.400) (1814,0.400) (1816,0.600) (1818,0.400) (1820,0.400) (1822,0.400) (1824,0.200) (1826,0.400) (1828,0.400) (1830,0.400) (1832,0.200) (1834,0.400) (1836,0.400) (1838,0.400) (1840,0.400) (1842,0.400) (1844,0.400) (1846,0.400) (1848,0.200) (1850,0.400) (1852,0.600) (1854,0.400) (1856,0.400) (1858,0.400) (1860,0.400) (1862,0.400) (1864,0.200) (1866,0.400) (1868,0.400) (1870,0.400) (1872,0.400) (1874,0.400) (1876,0.400) (1878,0.200) (1880,0.400) (1882,0.200) (1884,0.400) (1886,0.600) (1888,0.400) (1890,0.400) (1892,0.400) (1894,0.400) (1896,0.600) (1898,0.400) (1900,0.400) (1902,0.400) (1904,0.400) (1906,0.200) (1908,0.400) (1910,0.400) (1912,0.600) (1914,0.600) (1916,0.800) (1918,0.800) (1920,0.800) (1922,0.800) (1924,0.800) (1926,0.600) (1928,0.400) (1930,0.400) (1932,0.200) (1934,0.400) (1936,0.200) (1938,0.200) (1940,0.200) (1942,0.400) (1944,0.400) (1946,0.600) (1948,0.400) (1950,0.600) (1952,0.400) (1954,0.400) (1956,0.600) (1958,0.600) (1960,0.600) (1962,0.800) (1964,0.600) (1966,0.600) (1968,0.600) (1970,0.600) (1972,0.400) (1974,0.600) (1976,0.800) (1978,0.800) (1980,0.600) (1982,0.800) (1984,0.600) (1986,0.600) (1988,0.400) (1990,0.600) (1992,0.400) (1994,0.600) (1996,0.600) (1998,0.800) (2000,0.800)};
  \node[font=\sffamily\footnotesize, anchor=west] at (axis cs:2120,0.55) {$N=10$};
\end{axis}
\end{tikzpicture}\par\vspace{3mm}
\begin{tikzpicture}
\begin{axis}[bkplot, scale only axis, width=96mm, height=15mm, xmin=0, xmax=2080, ymin=0, ymax=1.15, xtick={500,1000,1500,2000}, ytick={0,0.5,1}, yticklabel style={/pgf/number format/fixed}, xticklabels={}, ylabel={$L/N$}]
  \addplot[bkOrange, line width=1pt] coordinates {(0,1.000) (20,0.840) (40,0.680) (60,0.560) (80,0.560) (100,0.560) (120,0.540) (140,0.520) (160,0.520) (180,0.520) (200,0.520) (220,0.480) (240,0.500) (260,0.440) (280,0.380) (300,0.400) (320,0.340) (340,0.380) (360,0.360) (380,0.380) (400,0.460) (420,0.440) (440,0.440) (460,0.420) (480,0.460) (500,0.500) (520,0.460) (540,0.440) (560,0.480) (580,0.480) (600,0.480) (620,0.520) (640,0.560) (660,0.560) (680,0.520) (700,0.440) (720,0.540) (740,0.440) (760,0.440) (780,0.420) (800,0.420) (820,0.420) (840,0.480) (860,0.460) (880,0.460) (900,0.440) (920,0.500) (940,0.480) (960,0.440) (980,0.420) (1000,0.440) (1020,0.480) (1040,0.520) (1060,0.540) (1080,0.500) (1100,0.440) (1120,0.420) (1140,0.520) (1160,0.520) (1180,0.480) (1200,0.480) (1220,0.480) (1240,0.500) (1260,0.480) (1280,0.500) (1300,0.520) (1320,0.500) (1340,0.460) (1360,0.460) (1380,0.480) (1400,0.440) (1420,0.540) (1440,0.600) (1460,0.640) (1480,0.580) (1500,0.560) (1520,0.540) (1540,0.580) (1560,0.580) (1580,0.520) (1600,0.480) (1620,0.520) (1640,0.540) (1660,0.540) (1680,0.620) (1700,0.560) (1720,0.500) (1740,0.520) (1760,0.540) (1780,0.480) (1800,0.500) (1820,0.500) (1840,0.460) (1860,0.480) (1880,0.500) (1900,0.520) (1920,0.500) (1940,0.580) (1960,0.580) (1980,0.600) (2000,0.580)};
  \node[font=\sffamily\footnotesize, anchor=west] at (axis cs:2120,0.55) {$N=100$};
\end{axis}
\end{tikzpicture}\par\vspace{3mm}
\begin{tikzpicture}
\begin{axis}[bkplot, scale only axis, width=96mm, height=15mm, xmin=0, xmax=2080, ymin=0, ymax=1.15, xtick={500,1000,1500,2000}, ytick={0,0.5,1}, yticklabel style={/pgf/number format/fixed}, xlabel={回}, ylabel={$L/N$}]
  \addplot[bkBlue, line width=1.3pt] coordinates {(0,1.000) (20,0.980) (40,0.960) (60,0.942) (80,0.924) (100,0.904) (120,0.890) (140,0.874) (160,0.858) (180,0.848) (200,0.832) (220,0.820) (240,0.812) (260,0.798) (280,0.788) (300,0.770) (320,0.764) (340,0.750) (360,0.734) (380,0.722) (400,0.708) (420,0.704) (440,0.696) (460,0.686) (480,0.676) (500,0.672) (520,0.670) (540,0.668) (560,0.668) (580,0.660) (600,0.658) (620,0.654) (640,0.654) (660,0.646) (680,0.642) (700,0.634) (720,0.622) (740,0.616) (760,0.612) (780,0.606) (800,0.602) (820,0.602) (840,0.596) (860,0.594) (880,0.600) (900,0.598) (920,0.596) (940,0.596) (960,0.588) (980,0.594) (1000,0.588) (1020,0.584) (1040,0.586) (1060,0.582) (1080,0.584) (1100,0.572) (1120,0.566) (1140,0.560) (1160,0.554) (1180,0.546) (1200,0.548) (1220,0.546) (1240,0.546) (1260,0.542) (1280,0.538) (1300,0.538) (1320,0.532) (1340,0.530) (1360,0.526) (1380,0.524) (1400,0.512) (1420,0.510) (1440,0.506) (1460,0.508) (1480,0.514) (1500,0.512) (1520,0.516) (1540,0.512) (1560,0.508) (1580,0.512) (1600,0.506) (1620,0.508) (1640,0.498) (1660,0.496) (1680,0.498) (1700,0.498) (1720,0.492) (1740,0.488) (1760,0.478) (1780,0.476) (1800,0.480) (1820,0.480) (1840,0.484) (1860,0.488) (1880,0.478) (1900,0.482) (1920,0.486) (1940,0.490) (1960,0.486) (1980,0.486) (2000,0.492)};
  \node[font=\sffamily\footnotesize, anchor=west] at (axis cs:2120,0.55) {$N=1000$};
\end{axis}
\end{tikzpicture}
\bkfigcaption{エーレンフェストのつぼ（たて軸は左の割合 $L/N$）。はじめは全部左（$L/N=1$）。玉が多いほど、ゆっくり半分に近づき、近づいたあとのゆれは小さい。$N=10$ は2回ごと、ほかは20回ごとに点をとった。}\label{ph:熱-15:fig-tsubo}
\end{bkfigure}

\phboxhead{グラフから読みとれること}
\begin{enumerate}[label={\bkcircled{\arabic*}}, leftmargin=2.4\zw, itemsep=1pt, topsep=2pt]
\item どの $N$ でも、左の割合は半分に近づきます。近づくまでの回数は、$N=100$ で百回ほど、$N=1000$ で千回ほどで、$N$ に比例します。
\item 近づいたあとのゆれは、$N=10$ では0から1まで大きく、$N=100$ では0.33〜0.64、$N=1000$ では0.48〜0.52の細い帯に入っています。
\item $N=10$ では、全部が片方に集まること（$L/N=1$ か0）も起こります。この例では、\num{1188}回めと\num{1610}回めに、全部が右に集まりました。$N=1000$ では、集まりそうにありません。
\end{enumerate}

\phboxhead{平均すると、どう近づく？}
半分からのずれを $d=L-N/2$ とします。1回動かすと、$L$ は確率 $L/N$ で1減り、確率 $(N-L)/N$ で1増えるので、平均の変化は
\[(-1)\times\frac{L}{N}+(+1)\times\frac{N-L}{N}=\frac{N-2L}{N}=-\frac{2}{N}\,d\]
です。ずれは、平均すると1回ごとに $(1-2/N)$ 倍になります。$t$ 回では $(1-2/N)^t$ 倍です。$N=1000$ で $t=500$ 回なら $0.998^{500}\approx0.37$（ほぼ $e^{-1}$。5-14②）で、平均の割合は $0.5+0.5\times0.37\approx0.68$ になります。図\ref{ph:熱-15:fig-tsubo}の500回めも約0.67です。ずれが $e^{-1}$ 倍になるまでの回数が $N/2$ なので、近づく回数は $N$ に比例します（読みとれること①）。

\phboxhead{プログラム：近づいたあとのゆれの大きさ}
ゆれの大きさを、標準偏差（3-14）ではかります。近づいたあとの $L$ を集めて、平均と標準偏差を出すプログラムです。
\begin{phcode}[Python]
import random
random.seed(1)
for N in [100, 10000]:
    L = N                      \# はじめは全部左
    data = []
    for step in range(200 * N):
        if random.random() < L / N:   \# 左の玉を選んだ
            L = L - 1
        else:
            L = L + 1
        if step >= 20 * N:            \# 近づいてから集める
            data.append(L)
    ave = sum(data) / len(data)
    sd = (sum((x - ave)**2 for x in data) / len(data)) ** 0.5
    print(N, round(ave, 1), round(sd, 1))
\end{phcode}
\begin{phcodeout}
100 50.1 5.1
10000 5002.3 49.4
\end{phcodeout}
平均はほぼ $N/2$、標準偏差は $N=100$ で約5、$N=\num{10000}$ で約50です。ゆれの大きさは、およそ $\sqrt N/2$ 個（$\sqrt{100}/2=5$、$\sqrt{\num{10000}}/2=50$）です。これはコインを $N$ 枚投げたときの、表の枚数のちらばりと同じです（3-17⑤の正規分布）。

\begin{phbaisuru}[100倍]{玉の数 $N$}
$N=100$ から $N=\num{10000}$ にすると、ゆれの個数は $\sqrt{100}=10$ 倍（5個→50個）になりますが、全体に対する割合は10分の1（5\,\%→0.5\,\%）になります。近づくまでの回数は100倍です。1 mol（$6\times10^{23}$ 個）なら、ゆれの割合は約 $6\times10^{-13}$ で、どんな測りかたでも見えません。だから、たくさんの分子でできた物は、ゆれのない、決まった温度・圧力をもつように見えるのです。
\end{phbaisuru}

\begin{seikaku}
全部が左にもどることは、つぼのモデルでも禁じられてはいません。全部が左の状態から始めて、もう一度全部が左にもどるまでの回数の平均は、ちょうど $2^N$ 回であることが知られています（カッツ、1947年）。$N=10$ なら約\num{1000}回なので、図\ref{ph:熱-15:fig-tsubo}の\num{2000}回の中で起こってもおかしくありません。$N=100$ では $1.3\times10^{30}$ 回で、1秒に1回動かしても、宇宙の年齢の約 $3\times10^{12}$ 倍かかります。
\end{seikaku}

%% ---------------------------------------------------------------
\bksubsection{4}{ボルツマンの関係：$S=k_{\mathrm B}\ln\Omega$}
\label{ph:熱-15:sec4}
①〜③で、変化は、状態の数 $\Omega$ が大きいマクロな状態へ進むとわかりました。一方、熱-6 では、まわりと熱のやりとりのない物の変化は、エントロピー $S$ が増える向きに進みました。$\Omega$ と $S$ は、どちらも「変化の向き」を表しています。2つの間には、どんな関係があるのでしょう。

\phboxhead{2つの物をあわせると}
\noindent\begin{minipage}[c]{0.56\linewidth}\raggedright
物 A のミクロな状態の数を $\Omega_{\mathrm A}$、物 B の数を $\Omega_{\mathrm B}$ とします。A の1つの状態ごとに、B は $\Omega_{\mathrm B}$ 通りの状態をとれるので、あわせた物の状態の数は、かけ算 $\Omega_{\mathrm A}\times\Omega_{\mathrm B}$ です（3-17②の積の法則）。ところが、エントロピーは、たし算 $S_{\mathrm A}+S_{\mathrm B}$ です（熱-6：熱÷温度は、それぞれの分をたす）。
\end{minipage}\hfill
\begin{minipage}[c]{0.4\linewidth}\centering
\begin{tikzpicture}[x=1mm, y=1mm]
  \node[bk node, minimum width=18mm, minimum height=10mm, font=\sffamily\footnotesize] (a) at (0,0) {A：3通り};
  \node[bk node, minimum width=18mm, minimum height=10mm, font=\sffamily\footnotesize] (b) at (26,0) {B：4通り};
  \node[bk node, minimum width=44mm, minimum height=10mm, font=\sffamily\footnotesize, draw=bkRed] (ab) at (13,-20) {あわせて $3\times4=12$ 通り};
  \draw[bk link] (a.south) -- (ab.north -| a.south);
  \draw[bk link] (b.south) -- (ab.north -| b.south);
  \node[qxp lab, anchor=north] at (13,-26.5) {$\ln12=\ln3+\ln4$};
\end{tikzpicture}
\end{minipage}

\par\vspace{4pt}
かけ算をたし算に変えるのは、対数です（3-12③：$\log(ab)=\log a+\log b$）。そこで、ボルツマンは、エントロピーを状態の数の対数で表しました。
\pheq{S=k_{\mathrm B}\ln\Omega}{J/K ＝ J/K × 1}
\label{ph:熱-15:boltzmann}これを\textbf{ボルツマンの関係}といいます。\bkidx{ぼるつまんのかんけい}{ボルツマンの関係}$k_{\mathrm B}$ はボルツマン定数（熱-14）で、ただの数 $\ln\Omega$ を J/K の単位に直す役目をします。ボルツマンが1877年の論文で、エントロピーと確率の関係を示し、この形の式と定数 $k$ は、プランクが1900〜1901年に書きました。$S$ が $\Omega$ とともに増えるので、「$\Omega$ が増える向きに変わる」ことと「$S$ が増える向きに変わる」ことは同じになります。

\begin{phyomi}[$S=k_{\mathrm B}\ln\Omega$]
「状態の数 $\Omega$ が何けたの数かを数えると、エントロピーになる」。$\ln\Omega$ は、$\Omega$ のけたの数のおよそ2.3倍です（$\ln10\approx2.3$）。並び方が多いほど、つまり、ミクロに見て「どの状態にいるかわからない」度合いが大きいほど、エントロピーは大きくなります。
\end{phyomi}

\begin{phbaisuru}{状態の数 $\Omega$}
$S$ は2倍にはなりません。$k_{\mathrm B}\ln2\approx\qty{9.6e-24}{J/K}$ だけ増えるだけです。$S$ を2倍にするには、$\Omega$ を2乗しなければなりません（$\ln\Omega^2=2\ln\Omega$）。2つの同じ物をあわせると $\Omega$ は2乗、$S$ は2倍です。
\end{phbaisuru}

\begin{phmijika}[1 mol の窒素の状態の数]
温度25\,℃、圧力\qty{100}{kPa}の窒素1 mol のエントロピーは、表の値で約\qty{192}{J/K}です。$\ln\Omega=S/k_{\mathrm B}\approx1.39\times10^{25}$ なので、$\Omega$ は10を約 $6\times10^{24}$ 回かけた数、$10^{6\times10^{24}}$ ほどです。トランプの $10^{68}$ とはくらべものにならない大きさです。
\end{phmijika}

\phboxhead{熱-6 の $R\ln2$ を、数えて求める}
熱-6 ④ では、理想気体1 mol が、仕切りを外して真空の側へ広がり、体積が2倍になると（自由膨張）、エントロピーが $R\ln2$ 増えることを、温度を一定に保ってゆっくり広げる道（可逆の道）で $\int\delta Q/T$ を計算して求めました。これを、数えて求めます。

箱の中を、同じ大きさの小さな部屋に区切って考えます。体積が2倍になると、1個の分子が入れる部屋の数は2倍になります。分子は互いにじゃまをしない（理想気体）ので、$N$ 個では、並び方の数が $2^N$ 倍になります。
\[\Delta S=k_{\mathrm B}\ln(2^N\,\Omega)-k_{\mathrm B}\ln\Omega=k_{\mathrm B}\ln2^N=Nk_{\mathrm B}\ln2\]
1 mol では $N=N_{\mathrm A}$ で、$N_{\mathrm A}k_{\mathrm B}=R$（気体定数。熱-3）なので、
\[\Delta S=R\ln2=8.314\times0.693\approx\qty{5.76}{J/K}\]
です（例題4）。熱を測って求めた熱-6 ④ の値と、並び方を数えた値が、ぴったり同じになりました。

\begin{phflow}[並び方の数からエントロピーの変化へ]
\phflowrow{$S=k_{\mathrm B}\ln\Omega$}{部屋の数は体積に比例\\分子 $N$ 個}{$\Delta S=Nk_{\mathrm B}\ln\dfrac{V_2}{V_1}$}
\phflowrow{$\Delta S=Nk_{\mathrm B}\ln\dfrac{V_2}{V_1}$}{1 mol（$N=N_{\mathrm A}$）\\$V_2=2V_1$}{$\Delta S=R\ln2\approx\qty{5.76}{J/K}$}
\end{phflow}

\phboxhead{温度は「エネルギーをもらったときの、$\Omega$ の増え方」}
次は、熱のやりとりを数えます。固体の原子は、決まった場所のまわりでゆれています。ここでは、原子のゆれのエネルギーが、決まった大きさ $\Delta\varepsilon$ の「\textbf{エネルギーの粒}」の何個分かで数えられるとします（なぜそうなるのかは⑧と、物理シリーズの『原子』（原-6）で）。ゆれる原子 $N$ 個に、粒 $j$ 個を分ける分け方の数は、粒 $j$ 個と仕切り $N-1$ 本を1列に並べる並べ方の数です（図\ref{ph:熱-15:fig-tsubu}）。
\pheq{\Omega(N,\ j)={}_{j+N-1}\mathrm{C}_j}{通り ＝ 通り}
\begin{bkfigure}
\begin{tikzpicture}[x=1mm, y=1mm]
  \foreach \row/\s/\a in {0/{1,1,0,0}/{2,0,0}, 1/{1,0,1,0}/{1,1,0}, 2/{1,0,0,1}/{1,0,1}, 3/{0,1,1,0}/{0,2,0}, 4/{0,1,0,1}/{0,1,1}, 5/{0,0,1,1}/{0,0,2}} {
    \begin{scope}[shift={(0,-\row*6.4)}]
      \foreach \v [count=\i] in \s {\ifnum\v=1 \fill[bkOrange] (\i*6,0) circle (1.7); \else \draw[bkInk, line width=1.4pt] (\i*6,-2.2) -- (\i*6,2.2); \fi}
      \node[qxp lab, anchor=west] at (31,0) {$\to$ $(\a)$};
    \end{scope}}
  \node[qxp lab, anchor=south west] at (29,3.2) {原子1・2・3の粒の数};
  \node[qxp lab, anchor=west, text=bkOrange!80!black] at (82,-6) {●：エネルギーの粒};
  \node[qxp lab, anchor=west] at (82,-12.4) {｜：仕切り（2本）};
  \node[qxp lab, anchor=west] at (82,-18.8) {${}_4\mathrm{C}_2=6$ 通り};
\end{tikzpicture}
\bkfigcaption{原子3個に粒2個を分ける分け方。粒2個と仕切り2本の並べ方で数えられる。仕切りの左が原子1、間が原子2、右が原子3の分。}\label{ph:熱-15:fig-tsubu}
\end{bkfigure}

原子3個の物 A と、原子3個の物 B を接触させ、全部で6個の粒を分け合うとします。A の粒の数 $j_{\mathrm A}$ ごとの、状態の数です。
\begin{bkcenter}
\begin{bktable}{colspec={Q[l,m,wd=38mm] *{7}{Q[c,m,wd=9mm]}}, rowsep=1.5pt, row{1}={bg=bkTheme!10}}
A の粒の数 $j_{\mathrm A}$ & 0 & 1 & 2 & 3 & 4 & 5 & 6 \\
B の粒の数 $j_{\mathrm B}$ & 6 & 5 & 4 & 3 & 2 & 1 & 0 \\
$\Omega_{\mathrm A}={}_{j_{\mathrm A}+2}\mathrm{C}_{j_{\mathrm A}}$ & 1 & 3 & 6 & 10 & 15 & 21 & 28 \\
$\Omega_{\mathrm B}$ & 28 & 21 & 15 & 10 & 6 & 3 & 1 \\
$\Omega_{\mathrm A}\times\Omega_{\mathrm B}$ & 28 & 63 & 90 & 100 & 90 & 63 & 28 \\
\end{bktable}
\bktabcaption{原子3個ずつの物 A・B で、粒6個を分け合う。全部で462通り（$28+63+\cdots+28$）。}\label{ph:熱-15:tab-kotai}
\end{bkcenter}
はじめに A が冷たく（$j_{\mathrm A}=0$）、B が熱かった（$j_{\mathrm B}=6$）とします。このマクロな状態のミクロな状態は28通りですが、粒をやりとりすると、もっと数の多い $j_{\mathrm A}=1$、2、3 のほうへ進みます。いちばん多いのは $j_{\mathrm A}=j_{\mathrm B}=3$（100通り）です。\textbf{熱が熱い物から冷たい物へ流れ、温度がそろうのは、そのほうが状態の数が多いから}なのです。原子の数を増やすと、半分ずつの分け方が、ますます飛びぬけて起こりやすくなります（図\ref{ph:熱-15:fig-kotai}）。
\begin{bkfigure}
\begin{tikzpicture}
\begin{axis}[bkplot, width=100mm, height=44mm, xmin=0, xmax=1.08, ymin=0, ymax=15.5, xtick={0.25,0.5,0.75,1}, ytick={5,10,15}, xlabel={$j_{\mathrm A}/j$}, ylabel={起こりやすさ}]
  \addplot[ycomb, bkGray!80!black, line width=3.2pt] coordinates {(0.0000,0.3636) (0.1667,0.8182) (0.3333,1.1688) (0.5000,1.2987) (0.6667,1.1688) (0.8333,0.8182) (1.0000,0.3636)};
  \addplot[bkOrange, line width=1.1pt, dashed] coordinates {(0.0000,0.0000) (0.0167,0.0000) (0.0333,0.0000) (0.0500,0.0000) (0.0667,0.0000) (0.0833,0.0000) (0.1000,0.0000) (0.1167,0.0001) (0.1333,0.0004) (0.1500,0.0011) (0.1667,0.0026) (0.1833,0.0060) (0.2000,0.0130) (0.2167,0.0261) (0.2333,0.0496) (0.2500,0.0893) (0.2667,0.1527) (0.2833,0.2491) (0.3000,0.3885) (0.3167,0.5806) (0.3333,0.8331) (0.3500,1.1499) (0.3667,1.5289) (0.3833,1.9604) (0.4000,2.4270) (0.4167,2.9035) (0.4333,3.3589) (0.4500,3.7597) (0.4667,4.0738) (0.4833,4.2741) (0.5000,4.3430) (0.5167,4.2741) (0.5333,4.0738) (0.5500,3.7597) (0.5667,3.3589) (0.5833,2.9035) (0.6000,2.4270) (0.6167,1.9604) (0.6333,1.5289) (0.6500,1.1499) (0.6667,0.8331) (0.6833,0.5806) (0.7000,0.3885) (0.7167,0.2491) (0.7333,0.1527) (0.7500,0.0893) (0.7667,0.0496) (0.7833,0.0261) (0.8000,0.0130) (0.8167,0.0060) (0.8333,0.0026) (0.8500,0.0011) (0.8667,0.0004) (0.8833,0.0001) (0.9000,0.0000) (0.9167,0.0000) (0.9333,0.0000) (0.9500,0.0000) (0.9667,0.0000) (0.9833,0.0000) (1.0000,0.0000)};
  \addplot[bkBlue, line width=1.3pt] coordinates {(0.0000,0.0000) (0.0083,0.0000) (0.0167,0.0000) (0.0250,0.0000) (0.0333,0.0000) (0.0417,0.0000) (0.0500,0.0000) (0.0583,0.0000) (0.0667,0.0000) (0.0750,0.0000) (0.0833,0.0000) (0.0917,0.0000) (0.1000,0.0000) (0.1083,0.0000) (0.1167,0.0000) (0.1250,0.0000) (0.1333,0.0000) (0.1417,0.0000) (0.1500,0.0000) (0.1583,0.0000) (0.1667,0.0000) (0.1750,0.0000) (0.1833,0.0000) (0.1917,0.0000) (0.2000,0.0000) (0.2083,0.0000) (0.2167,0.0000) (0.2250,0.0000) (0.2333,0.0000) (0.2417,0.0000) (0.2500,0.0000) (0.2583,0.0000) (0.2667,0.0000) (0.2750,0.0000) (0.2833,0.0000) (0.2917,0.0000) (0.3000,0.0000) (0.3083,0.0000) (0.3167,0.0000) (0.3250,0.0000) (0.3333,0.0000) (0.3417,0.0000) (0.3500,0.0000) (0.3583,0.0001) (0.3667,0.0003) (0.3750,0.0010) (0.3833,0.0036) (0.3917,0.0113) (0.4000,0.0327) (0.4083,0.0867) (0.4167,0.2106) (0.4250,0.4687) (0.4333,0.9570) (0.4417,1.7936) (0.4500,3.0875) (0.4583,4.8844) (0.4667,7.1045) (0.4750,9.5045) (0.4833,11.6984) (0.4917,13.2497) (0.5000,13.8111) (0.5083,13.2497) (0.5167,11.6984) (0.5250,9.5045) (0.5333,7.1045) (0.5417,4.8844) (0.5500,3.0875) (0.5583,1.7936) (0.5667,0.9570) (0.5750,0.4687) (0.5833,0.2106) (0.5917,0.0867) (0.6000,0.0327) (0.6083,0.0113) (0.6167,0.0036) (0.6250,0.0010) (0.6333,0.0003) (0.6417,0.0001) (0.6500,0.0000) (0.6583,0.0000) (0.6667,0.0000) (0.6750,0.0000) (0.6833,0.0000) (0.6917,0.0000) (0.7000,0.0000) (0.7083,0.0000) (0.7167,0.0000) (0.7250,0.0000) (0.7333,0.0000) (0.7417,0.0000) (0.7500,0.0000) (0.7583,0.0000) (0.7667,0.0000) (0.7750,0.0000) (0.7833,0.0000) (0.7917,0.0000) (0.8000,0.0000) (0.8083,0.0000) (0.8167,0.0000) (0.8250,0.0000) (0.8333,0.0000) (0.8417,0.0000) (0.8500,0.0000) (0.8583,0.0000) (0.8667,0.0000) (0.8750,0.0000) (0.8833,0.0000) (0.8917,0.0000) (0.9000,0.0000) (0.9083,0.0000) (0.9167,0.0000) (0.9250,0.0000) (0.9333,0.0000) (0.9417,0.0000) (0.9500,0.0000) (0.9583,0.0000) (0.9667,0.0000) (0.9750,0.0000) (0.9833,0.0000) (0.9917,0.0000) (1.0000,0.0000)};
  \node[font=\sffamily\footnotesize, text=bkGray!70!black, anchor=south west] at (axis cs:0.78,1.4) {原子3個ずつ（棒）};
  \node[font=\sffamily\footnotesize, text=bkOrange!75!black, anchor=west] at (axis cs:0.62,4.8) {30個ずつ（破線）};
  \node[font=\sffamily\footnotesize, text=bkBlue!80!black, anchor=west] at (axis cs:0.56,12.5) {300個ずつ（実線）};
\end{axis}
\end{tikzpicture}
\bkfigcaption{A と B が同じ大きさのとき、A がもつ粒の割合 $j_{\mathrm A}/j$ の起こりやすさ（粒の数 $j$ は原子の数の2倍。たて軸は確率に $j$ をかけて、曲線の下の面積が1になるようにした）。原子が多いほど、半分ずつの所に鋭く集まる。}\label{ph:熱-15:fig-kotai}
\end{bkfigure}


では、粒がどこまで移ると止まるのでしょう。A が粒を1個もらうと $\ln\Omega_{\mathrm A}$ が増え、B は1個失って $\ln\Omega_{\mathrm B}$ が減ります。原子3個の物が粒を1個もらうときの $\ln\Omega$ の増え方を計算すると、
\begin{bkcenter}
\begin{bktable}{colspec={Q[l,m,wd=40mm] *{6}{Q[c,m,wd=11mm]}}, rowsep=1.5pt, row{1}={bg=bkTheme!10}}
粒の数 $j$ → $j+1$ & 0→1 & 1→2 & 2→3 & 3→4 & 4→5 & 5→6 \\
$\ln\Omega$ の増え & 1.10 & 0.69 & 0.51 & 0.41 & 0.34 & 0.29 \\
\end{bktable}
\bktabcaption{粒が少ない（冷たい）物ほど、1個もらったときの $\ln\Omega$ の増え方が大きい。}\label{ph:熱-15:tab-fue}
\end{bkcenter}
となります。粒の少ない物（冷たい物）ほど、1個もらったときの増え方が大きいのです。A が冷たいうちは、A の増え方が B の減り方より大きいので、粒が A へ移ると全体の $\Omega$ が増えます。2つの増え方がつり合った所（3個と3個）で、流れは止まります。そこで、「エネルギーを1つもらったときの $S=k_{\mathrm B}\ln\Omega$ の増え方」で温度を決めます。
\pheq{\frac{1}{T}=\frac{\Delta S}{\Delta U}\quad(\text{体積を変えないとき})}{1/K ＝ (J/K)÷J}
$\Delta U$ は、もらったエネルギー（内部エネルギーの増え）です。増え方の大きい冷たい物ほど、$1/T$ が大きく、$T$ が小さくなっています。仕事をしないで熱 $Q$ をもらうなら $\Delta U=Q$ なので、この式は熱-6 の $\Delta S=Q/T$ そのものです。\textbf{数えて求めたエントロピーと、熱を測って求めたエントロピーは、同じものだった}のです。

\phboxhead{プログラム：もっと大きな物でたしかめる}
原子300個の物 A と200個の物 B が、粒100個を分け合うとき、いちばん多い分け方と、A の粒が50〜70個になる確率、全部の粒が B にある確率を出します。
\begin{phcode}[Python]
from math import comb
nA, nB, j = 300, 200, 100      \# A と B の原子の数、粒の数
def omega(N, j):               \# N 個に j 粒を分ける数
    return comb(j + N - 1, j)
tot = comb(j + nA + nB - 1, j) \# 全部の分け方
def both(a):                   \# A が a 粒のときの数
    return omega(nA, a) * omega(nB, j - a)
best = max(range(j + 1), key=both)
print(best)
p = sum(both(a) for a in range(50, 71)) / tot
print(round(p, 3))
print(both(0) / tot)
\end{phcode}
\begin{phcodeout}
60
0.95
2.9931326169281273e-35
\end{phcodeout}
いちばん多いのは、A が60個、B が40個のときで、原子1個あたりの粒の数がどちらも0.2個にそろっています（温度がそろう）。A が50〜70個になる確率は0.95、全部の粒が B に残る確率は $3\times10^{-35}$ です。原子が $10^{23}$ 個ある本物の固体なら、温度がそろったあと、そこからずれることは、事実上ありません。

\begin{seikaku}
等確率の原理は、⑤から先の式でも使う仮定です。どんな物でも成り立つことは、まだ証明されていません（⑨のあとの「まだわかっていないこと」）。また、$S=k_{\mathrm B}\ln\Omega$ は、エントロピーの「変化」だけでなく、値そのものも決めます。$\Omega=1$（状態がただ1つ）なら $S=0$ です。温度を0 K に近づけると、多くの物はいちばんエネルギーの低い1つの状態に落ち着くので、エントロピーは0に近づきます（熱力学の第三法則）。
\end{seikaku}

\phkokomade

\Needspace{16\baselineskip}\phstage{shiki}{起こりやすさを式で導く}{5}

\bksubsection{5}{大きな数の階乗：スターリングの近似}
\label{ph:熱-15:sec5}
分子は $10^{23}$ 個もあるので、⑥では $N!$ のような大きな数の対数が要ります。電卓では、$70!$ ほどで計算できなくなります。そこで、$\ln n!$ を近似する式を作ります。対数はかけ算をたし算にするので、
\[\ln n!=\ln(1\times2\times\cdots\times n)=\ln1+\ln2+\ln3+\cdots+\ln n\]
です。これは、はば1、高さ $\ln k$ の長方形の面積を、$k=2$ から $n$ までたしたものと同じです（図\ref{ph:熱-15:fig-stirling}。$\ln1=0$）。5-10と同じように、細い長方形の面積の和は、曲線 $y=\ln x$ の下の面積（積分）とほとんど同じです。そこで、和を積分でおきかえます。
\begin{bkfigure}
\begin{tikzpicture}
\begin{axis}[bkplot, width=112mm, height=46mm, xmin=0, xmax=11.2, ymin=0, ymax=2.7, xtick={1,2,3,4,5,6,7,8,9,10}, ytick={1,2}, xlabel={$x$}, ylabel={$y$}]
  \foreach \k in {2,...,10} {\edef\tmp{\noexpand\fill[bkOrange!25] (axis cs:\k-1,0) rectangle (axis cs:\k,{ln(\k)});
    \noexpand\draw[bkOrange!80!black, line width=0.5pt] (axis cs:\k-1,0) rectangle (axis cs:\k,{ln(\k)});}\tmp}
  \addplot[bkcurve, domain=1:10.6, samples=60] {ln(x)};
  \node[font=\sffamily\footnotesize, text=bkBlue!80!black, anchor=south west] at (axis cs:6.3,2.2) {$y=\ln x$};
  \node[font=\sffamily\footnotesize, text=bkOrange!75!black, fill=white, inner sep=1pt] at (axis cs:5.5,0.75) {長方形の面積の和 $=\ln10!$};
\end{axis}
\end{tikzpicture}
\bkfigcaption{$\ln10!=\ln2+\ln3+\cdots+\ln10$ は、長方形の面積の和。曲線 $y=\ln x$ の下の面積（1から10まで）より少し大きい。}\label{ph:熱-15:fig-stirling}
\end{bkfigure}

\phboxhead{$\ln x$ の積分}
$\ln x$ を微分すると $1/x$ でした（5-14⑦）。では、微分すると $\ln x$ になる関数は何でしょう。$F(x)=x\ln x-x$ を、積の微分（5-6②）で微分してみます。
\[F'(x)=\underbrace{1\times\ln x+x\times\frac{1}{x}}_{x\ln x\ \text{の微分}}-1=\ln x+1-1=\ln x\]
$F(x)$ は $\ln x$ の原始関数です。定積分は $F(n)-F(1)$（5-11④）なので、
\[\int_1^n\ln x\,\dd x=(n\ln n-n)-(1\times\ln1-1)=n\ln n-n+1\]
$n$ が大きいと、最後の1は $n\ln n$ にくらべて小さいので、すてます。
\pheq{\ln n!\approx n\ln n-n}{1 ＝ 1}
\label{ph:熱-15:stirling}これを\textbf{スターリングの近似}といいます。\bkidx{すたーりんぐのきんじ}{スターリングの近似}
\begin{bkcenter}
\begin{bktable}{colspec={Q[c,m,wd=14mm] Q[c,m,wd=26mm] Q[c,m,wd=26mm] Q[c,m,wd=22mm]}, rowsep=1.5pt, row{1}={bg=bkTheme!10}}
$n$ & $\ln n!$ & $n\ln n-n$ & ずれ（\%） \\
5 & 4.79 & 3.05 & 36 \\
10 & 15.10 & 13.03 & 14 \\
52 & 156.36 & 153.46 & 1.9 \\
100 & 363.74 & 360.52 & 0.89 \\
\num{1000} & \num{5912.13} & \num{5907.76} & 0.074 \\
\end{bktable}
\bktabcaption{スターリングの近似は、$n$ が大きいほど当たる（ずれ＝差÷$\ln n!$）。}\label{ph:熱-15:tab-stirling}
\end{bkcenter}

\begin{phyomi}[$\ln n!\approx n\ln n-n$]
両辺を $e$ の指数にもどすと、
\[n!\approx n^ne^{-n}=\Bigl(\frac{n}{e}\Bigr)^n\]
「$n!$ は、$1$ から $n$ までの $n$ 個の数を、全部 $n/e$（$n$ の約0.37倍。まん中の $n/2$ より少し小さい）にそろえてかけたくらい」。
\end{phyomi}

\begin{phbaisuru}[10倍]{$n$}
$n=100$ から\num{1000}にすると、$\ln n!$ は 363.7 から\num{5912.1}へ、約16倍になります（$n\ln n$ の $\ln n$ も大きくなるから）。ずれの割合は 0.89\,\% から 0.074\,\% へ、約12分の1です。分子の数（$10^{23}$）では、ずれはまったく気になりません。
\end{phbaisuru}

\begin{seikaku}
もっと正確な形は $\ln n!\approx n\ln n-n+\tfrac{1}{2}\ln(2\pi n)$ です。$n=10$ で15.096（本当は15.104）、$52!$ では $8.05\times10^{67}$（本当は $8.07\times10^{67}$）になります。$\tfrac{1}{2}\ln(2\pi n)$ の部分は、$n$ が大きいと $n\ln n$ にくらべて小さいので、この章ではすてます。
\end{seikaku}

\phboxhead{使ってみる：半分ずつの並び方の数}
②の表の「半分ずつの並び方」の対数を、スターリングの近似で計算します。
\[{}_N\mathrm{C}_{N/2}=\frac{N!}{(N/2)!\,(N/2)!}\]
なので、
\[\ln{}_N\mathrm{C}_{N/2}\approx(N\ln N-N)-2\Bigl(\frac{N}{2}\ln\frac{N}{2}-\frac{N}{2}\Bigr)=N\ln N-N\ln\frac{N}{2}=N\ln2\]
つまり ${}_N\mathrm{C}_{N/2}$ は、対数で見ると全部の並び方 $2^N$ とほとんど同じです（$N=1000$ で、$\ln$ の値の比は0.995）。②で見た「ほとんど全部の並び方が、半分のあたりにある」ことを、式で確かめられました。

%% ---------------------------------------------------------------
\bksubsection{6}{いちばん起こりやすい分け方：ボルツマン分布}
\label{ph:熱-15:sec6}
分子1個がとれるエネルギーを、とびとびの値 $\varepsilon_0=0<\varepsilon_1<\varepsilon_2<\cdots$ とします。高さのちがう「棚」に、$N$ 個の分子を置いていくと考えます（図\ref{ph:熱-15:fig-tana}）。棚 $i$ にいる分子の数を $n_i$ とすると、全体の数と全体のエネルギーが決まっています。
\[n_0+n_1+n_2+\cdots=N,\qquad \varepsilon_0n_0+\varepsilon_1n_1+\varepsilon_2n_2+\cdots=U\]
この2つの条件のもとで、どの分け方 $(n_0,\ n_1,\ n_2,\ \ldots)$ がいちばん起こりやすいでしょう。④の固体の粒の分け方と同じ問いを、一般の形で考えます。

\noindent\begin{minipage}[c]{0.5\linewidth}\centering
\begin{tikzpicture}[x=1mm, y=1mm]
  \foreach \y/\t in {0/{$\varepsilon_0=0$},10/{$\varepsilon_1$},20/{$\varepsilon_2$},30/{$\varepsilon_3$}} {\draw[bkInk, line width=0.9pt] (14,\y) -- (56,\y); \node[qxp lab, anchor=east] at (12.5,\y+0.3) {\t};}
  \foreach \x/\c/\l in {18/bkBlue/A, 24/bkRed/B, 30/bkGreen/C, 36/bkOrange/D} {\qxpm{(\x,2.2)}{\c}{\l}}
  \foreach \x/\c/\l in {22/bkPurple/E, 28/bkBlue/F} {\qxpm{(\x,12.2)}{\c}{\l}}
  \qxpm{(25,22.2)}{bkRed}{G}
  \node[qxp lab, anchor=west] at (44,4) {$n_0=4$};
  \node[qxp lab, anchor=west] at (44,14) {$n_1=2$};
  \node[qxp lab, anchor=west] at (44,24) {$n_2=1$};
  \node[qxp lab, anchor=west] at (44,34) {$n_3=0$};
\end{tikzpicture}
\end{minipage}\hfill
\begin{minipage}[c]{0.47\linewidth}\raggedright
\bkfigcaption{エネルギーの棚に分子を置く。分子 A〜G の7個。}\label{ph:熱-15:fig-tana}
\end{minipage}

\phboxhead{分け方の数}
分子に番号をつけて1列に並べ（$N!$ 通り）、はじめの $n_0$ 個を棚0に、次の $n_1$ 個を棚1に、…と置きます。同じ棚の中で順番を入れかえても、置き方は同じなので、その分（$n_0!$、$n_1!$、…）でわります。
\[\Omega=\frac{N!}{n_0!\,n_1!\,n_2!\cdots}\]
図\ref{ph:熱-15:fig-tana}の7個なら
\[\frac{7!}{4!\,2!\,1!}=\frac{5040}{48}=105\ \text{通り}\]
です。どの $n_i$ も大きいとして、スターリングの近似を使うと、
\[\ln\Omega\approx\ln N!-\sum_i\bigl(n_i\ln n_i-n_i\bigr)\]
です（$\sum_i$ は、棚 $i=0,\ 1,\ 2,\ \ldots$ の分を全部たす記号。5-10③）。$n_i$ はとても大きいので、なめらかに変えられる数として微分します。$n_i\ln n_i-n_i$ を $n_i$ で微分すると $\ln n_i+1-1=\ln n_i$ なので、
\[\frac{\partial\ln\Omega}{\partial n_i}=-\ln n_i\]

\phboxhead{条件が2つのときの未定乗数}
8-4のもっと知りたいでは、道（条件）の上でいちばん高い所では、$\mathrm{grad}$ どうしが平行になりました。条件が2つのときは、2つの条件を両方とも変えない動かし方では、$\ln\Omega$ が変わってはいけません。そのような動かし方は、2つの条件の $\mathrm{grad}$ のどちらとも直角です。だから、いちばん大きい所では、$\mathrm{grad}\,\ln\Omega$ は、2つの条件の $\mathrm{grad}$ をそれぞれ何倍かしてたしたものになります。その倍の数（未定乗数）を $\alpha$ と $\beta$ とします（8-4では1つなので $\mu$ と書きましたが、この章では $\mu$ を⑨で別の意味に使います）。1つめの条件 $\sum n_i$ を $n_i$ で微分すると1、2つめの条件 $\sum\varepsilon_in_i$ を微分すると $\varepsilon_i$ なので、
\[-\ln n_i=\alpha\times1+\beta\times\varepsilon_i\qquad\text{だから}\qquad n_i=e^{-\alpha}\,e^{-\beta\varepsilon_i}\]
棚の分子の数は、エネルギー $\varepsilon_i$ が大きいほど、$e^{-\beta\varepsilon_i}$ の形で減っていきます。全部たすと $N$ になるように $e^{-\alpha}$ を決めると、$n_i/N=e^{-\beta\varepsilon_i}\big/\sum_je^{-\beta\varepsilon_j}$ です。

\phboxhead{$\beta$ の正体は温度}
分け方をいちばん起こりやすい形にしたまま、$N$ を変えずに、全体のエネルギーを $\Delta U$ だけ増やします。$\ln\Omega$ の増えは、
\[\Delta\ln\Omega=\sum_i(-\ln n_i)\,\Delta n_i=\alpha\sum_i\Delta n_i+\beta\sum_i\varepsilon_i\,\Delta n_i=\alpha\times0+\beta\,\Delta U\]
です。$S=k_{\mathrm B}\ln\Omega$ なので $\Delta S=k_{\mathrm B}\beta\,\Delta U$。④の $1/T=\Delta S/\Delta U$ とくらべると、
\[\beta=\frac{1}{k_{\mathrm B}T}\]
です。まとめると、温度 $T$ で、分子がエネルギー $\varepsilon_i$ の状態にいる確率は
\pheq{P_i=\frac{e^{-\varepsilon_i/(k_{\mathrm B}T)}}{Z},\qquad Z=\sum_ie^{-\varepsilon_i/(k_{\mathrm B}T)}}{1 ＝ 1}
\label{ph:熱-15:bunpu}です。これを\textbf{ボルツマン分布}、$e^{-\varepsilon/(k_{\mathrm B}T)}$ を\textbf{ボルツマン因子}、$Z$ を\textbf{分配関数}といいます。\bkidx{ぼるつまんぶんぷ}{ボルツマン分布}\bkidx{ぶんぱいかんすう}{分配関数}8-4のもっと知りたいで「物理では、決まった全体のエネルギーの中で、いちばん起こりやすい分け方をさがすときに使います」と書いた未定乗数法の使い道が、これです。

\begin{phyomi}[$P_i\propto e^{-\varepsilon_i/(k_{\mathrm B}T)}$]
「エネルギーが $k_{\mathrm B}T$ 増えるごとに、そこにいる分子の数は $1/e$（約0.37）倍に減る」。$k_{\mathrm B}T$ は「温度 $T$ で分子が気軽に出し入れできるエネルギーの大きさ」で、300 K で $\qty{4.1e-21}{J}$ です。$k_{\mathrm B}T$ よりずっと高い棚には、ほとんど分子がいません。
\end{phyomi}

\begin{phbaisuru}{温度 $T$}
$\varepsilon/(k_{\mathrm B}T)$ が半分になるので、高い棚にも分子が上がります。もとの温度で $\varepsilon=2k_{\mathrm B}T$ の棚は、ボルツマン因子が $e^{-2}\approx0.14$ から $e^{-1}\approx0.37$ になります（図\ref{ph:熱-15:fig-bairitsu}）。低い棚の分子が減り、高い棚の分子が増えて、全体のエネルギーが増えます。
\par\smallskip\noindent\begin{tikzpicture}
\begin{axis}[bkplot, width=84mm, height=40mm, xmin=0, xmax=4.3, ymin=0, ymax=1.12, xtick={1,2,3,4}, ytick={0.14,0.37,1}, yticklabel style={/pgf/number format/fixed}, xlabel={$\varepsilon/(k_{\mathrm B}T_0)$}, xlabel style={at={(axis description cs:0.5,0)}, anchor=north, yshift=-4.5mm}, ylabel={ボルツマン因子}]
  \addplot[bkcurve, domain=0:4.1] {exp(-x)};
  \addplot[bkcurve2, dashed, domain=0:4.1] {exp(-x/2)};
  \fill[bkRed] (axis cs:2,0.1353) circle (1.8pt) (axis cs:2,0.3679) circle (1.8pt);
  \node[font=\sffamily\footnotesize, text=bkBlue!80!black, anchor=west] at (axis cs:4.2,0.01) {温度 $T_0$（実線）};
  \node[font=\sffamily\footnotesize, text=bkOrange!75!black, anchor=west] at (axis cs:4.2,0.25) {温度 $2T_0$（破線）};
\end{axis}
\end{tikzpicture}
\bkfigcaption{温度を2倍にすると、エネルギーの高い棚のボルツマン因子が大きくなる。}\label{ph:熱-15:fig-bairitsu}
\end{phbaisuru}

\begin{phmijika}[夏は食べ物がいたみやすい]
化学反応が起こるには、ぶつかった分子が、ある大きさ以上のエネルギー（反応の山）をもっている必要があります。そのような分子の割合はボルツマン因子に比例するので、反応の速さは温度で大きく変わります。山の高さが分子1個あたり\qty{0.5}{eV}（1 mol で約\qty{48}{kJ}。eV は⑦）なら、20\,℃から30\,℃にすると、$e^{-\varepsilon/(k_{\mathrm B}T)}$ は約1.9倍になります。反応の速さが温度でこの形に変わることは、実験の式（アレニウスの式、1889年）として知られていて、ボルツマン因子がその理由を説明しています。
\end{phmijika}

\begin{motto}[4]
\phboxhead{もう1つの道：大きな物とつながった、1個の分子}
ボルツマン分布は、④の考えだけからも導けます。1個の分子（小さな物）が、とても大きな物（まわり。温度 $T$）とエネルギーをやりとりしているとします。全体のエネルギーは $U_{\text{全}}$ で決まっています。分子がエネルギー $\varepsilon_i$ の状態にいるとき、まわりのエネルギーは $U_{\text{全}}-\varepsilon_i$ で、そのときの全体の状態の数は、まわりの状態の数 $\Omega_{\text{まわり}}(U_{\text{全}}-\varepsilon_i)$ です。まわりは大きいので、エネルギーを $\varepsilon_i$ だけ失っても温度は変わらず、④の式から、エントロピーは $\varepsilon_i/T$ だけ減ります。$\Omega=e^{S/k_{\mathrm B}}$ なので、
\[\Omega_{\text{まわり}}(U_{\text{全}}-\varepsilon_i)=e^{S_{\text{まわり}}/k_{\mathrm B}-\varepsilon_i/(k_{\mathrm B}T)}\ \propto\ e^{-\varepsilon_i/(k_{\mathrm B}T)}\]
等確率の原理から、確率は状態の数に比例するので、$P_i\propto e^{-\varepsilon_i/(k_{\mathrm B}T)}$ です。⑥と同じ形になりました。分子がたくさんある気体だけでなく、温度の決まったまわりとつながった、どんな小さな物にも使えることがわかります。
\end{motto}

\phboxhead{重力の中の分子：大気の式を導く}
地面からの高さ $z$ にある分子の位置エネルギーは $mgz$（力-4。$m$ は分子1個の質量）です。温度が一定なら、高さ $z$ の1 m$^3$ あたりの分子の数 $n_V(z)$ は、ボルツマン因子に比例します。
\begin{phflow}[ボルツマン分布から]
\phflowrow{$P_i\propto e^{-\varepsilon_i/(k_{\mathrm B}T)}$}{重力の中の分子\\$\varepsilon=mgz$・$T$ 一定}{$n_V(z)\propto e^{-mgz/(k_{\mathrm B}T)}$}
\phflowrow{$n_V(z)\propto e^{-mgz/(k_{\mathrm B}T)}$}{理想気体 $p=n_Vk_{\mathrm B}T$\\$H=k_{\mathrm B}T/(mg)$}{$p=p_0e^{-z/H}$\\（熱-3 ⑥ と同じ）}
\phflowrow{$P_i\propto e^{-\varepsilon_i/(k_{\mathrm B}T)}$}{$x$ 向きの運動\\$\varepsilon=\frac{1}{2}m{v_x}^2$}{$e^{-m{v_x}^2/(2k_{\mathrm B}T)}$ に比例\\（マクスウェル分布：熱-14）}
\end{phflow}
$H=k_{\mathrm B}T/(mg)$ は、分子の数が $1/e$ 倍になる高さです。$k_{\mathrm B}/m=R/M$（$M$ はモル質量）なので、$H=RT/(Mg)$ とも書けます。熱-3 ⑥ では、力のつり合い（$\dd p/\dd z=-\rho g$）と状態方程式から同じ式を作りました。窒素なら、15\,℃で $H\approx\qty{8.7}{km}$ です（例題6）。

%% ---------------------------------------------------------------
\bksubsection{7}{分配関数：2つの準位だけをもつ物}
\label{ph:熱-15:sec7}
\phboxhead{分配関数から平均のエネルギーを出す}
$\beta=1/(k_{\mathrm B}T)$ と書くと、$Z=\sum_ie^{-\beta\varepsilon_i}$ です。$Z$ を $\beta$ で微分すると、$e^{-\beta\varepsilon_i}$ の微分は $-\varepsilon_ie^{-\beta\varepsilon_i}$（5-14⑤）なので、
\[\frac{\dd Z}{\dd\beta}=-\sum_i\varepsilon_ie^{-\beta\varepsilon_i}\]
1個の分子の平均のエネルギーは、エネルギー×確率の和なので、
\pheq{\langle\varepsilon\rangle=\sum_i\varepsilon_iP_i=\frac{1}{Z}\sum_i\varepsilon_ie^{-\beta\varepsilon_i}=-\frac{1}{Z}\frac{\dd Z}{\dd\beta}=-\frac{\dd\ln Z}{\dd\beta}}{J ＝ J}
（最後は、合成関数の微分 $(\ln Z)'=Z'/Z$。5-6⑤・5-14⑦）。$N$ 個なら $U=N\langle\varepsilon\rangle$ です。\textbf{$Z$ さえ計算できれば、平均のエネルギーも、そこから比熱も出せる}ので、$Z$ は統計力学でいちばん大切な関数です。

\phboxhead{2準位の物}
エネルギーが0と $\Delta\varepsilon$ の2つの状態（\textbf{準位}）しかない粒子を考えます。磁場の中に置いた、小さな磁石のような原子の核や電子は、向きが2つしかとれず、このような2準位の物になります（なぜ2つなのかは、『原子』で）。
\[Z=e^{0}+e^{-\Delta\varepsilon/(k_{\mathrm B}T)}=1+e^{-\Delta\varepsilon/(k_{\mathrm B}T)}\]
上の準位と下の準位の数の比は、ボルツマン因子そのものです。
\pheq{\frac{\text{上の数}}{\text{下の数}}=e^{-\Delta\varepsilon/(k_{\mathrm B}T)}}{1 ＝ 1}
原子の世界のエネルギーは、\textbf{電子ボルト}（eV）で表すと便利です。\qty{1}{eV}は、電子1個が1 V の電圧で加速されて得るエネルギーで、\qty{1}{eV}＝\qty{1.602176634e-19}{J}です。300 K の $k_{\mathrm B}T$ は約\qty{0.0259}{eV}（ほぼ $\tfrac{1}{40}$ eV）です。差が\qty{0.1}{eV}なら、300 K で比は $e^{-0.1/0.0259}\approx0.021$ です（例題7）。

\phboxhead{2準位の物のエネルギーと比熱}
上にいる確率は、分母と分子に $e^{\Delta\varepsilon/(k_{\mathrm B}T)}$ をかけて
\[\frac{e^{-\Delta\varepsilon/(k_{\mathrm B}T)}}{1+e^{-\Delta\varepsilon/(k_{\mathrm B}T)}}=\frac{1}{e^{\Delta\varepsilon/(k_{\mathrm B}T)}+1}\]
です。$N$ 個のエネルギーは
\[U=\frac{N\Delta\varepsilon}{e^{\Delta\varepsilon/(k_{\mathrm B}T)}+1}\]
です。$x=\Delta\varepsilon/(k_{\mathrm B}T)$ とおいて、$U$ を温度で微分します。合成関数の微分（5-6⑤）で
\[\frac{\dd x}{\dd T}=-\frac{\Delta\varepsilon}{k_{\mathrm B}T^2}\]
を使うと、$N$ 個の熱容量は
\pheq{C_{\mathrm{th}}=\frac{\dd U}{\dd T}=Nk_{\mathrm B}\,\frac{x^2e^x}{(e^x+1)^2},\qquad x=\frac{\Delta\varepsilon}{k_{\mathrm B}T}}{J/K ＝ J/K}
となります。グラフにすると、熱容量はある温度で山になります（図\ref{ph:熱-15:fig-2jun}）。この山を\textbf{ショットキーの山}（ショットキー異常）といいます。
\begin{bkfigure}
\begin{tikzpicture}
\begin{axis}[bkplot, width=112mm, height=46mm, xmin=0, xmax=3.3, ymin=0, ymax=0.64, xtick={0.42,1,2,3}, ytick={0.1,0.2,0.3,0.4,0.5}, yticklabel style={/pgf/number format/fixed}, xlabel={$k_{\mathrm B}T/\Delta\varepsilon$}]
  \addplot[bkcurve, domain=0.05:3.2, samples=120] {(1/x)^2*exp(1/x)/(exp(1/x)+1)^2};
  \addplot[bkcurve2, dashed, domain=0.05:3.2, samples=120] {1/(exp(1/x)+1)};
  \addplot[bkthin, dotted, domain=0:3.2] {0.5};
  \fill[bkRed] (axis cs:0.417,0.439) circle (1.8pt);
  \node[font=\sffamily\footnotesize, text=bkRed!80!black, anchor=south west] at (axis cs:0.40,0.52) {山：$k_{\mathrm B}T\approx0.42\Delta\varepsilon$、$C_{\mathrm{th}}\approx0.44Nk_{\mathrm B}$};
  \node[font=\sffamily\footnotesize, text=bkBlue!80!black, anchor=south west] at (axis cs:2.0,0.095) {$C_{\mathrm{th}}/(Nk_{\mathrm B})$（実線）};
  \node[font=\sffamily\footnotesize, text=bkOrange!75!black, anchor=north west] at (axis cs:1.9,0.31) {$U/(N\Delta\varepsilon)$（破線）};
\end{axis}
\end{tikzpicture}
\bkfigcaption{2準位の物のエネルギー（破線）と熱容量（実線）。エネルギーは高温で $\tfrac{1}{2}N\Delta\varepsilon$（点線）に近づく。}\label{ph:熱-15:fig-2jun}
\end{bkfigure}
\begin{phyomi}[ショットキーの山]
「温度が低すぎると、粒子は上の準位に上がれず、熱をもらっても使い道がない。温度が高すぎると、上と下がもう半分ずつで、それ以上は上がれない。そのちょうど間（$k_{\mathrm B}T\approx0.42\Delta\varepsilon$）で、熱をいちばんよく吸いこむ」。
\end{phyomi}
\begin{phbaisuru}{準位の差 $\Delta\varepsilon$}
グラフの横軸が $k_{\mathrm B}T/\Delta\varepsilon$ なので、山の温度も2倍になります。山の高さ（$0.44Nk_{\mathrm B}$）は変わりません。差が\qty{0.1}{eV}なら山は約480 K、\qty{0.2}{eV}なら約970 K です。
\end{phbaisuru}
\begin{seikaku}
温度をどんなに上げても、上の準位の数は、下の準位の数をこえません（$e^{-\Delta\varepsilon/(k_{\mathrm B}T)}<1$）。レーザーでは、上の準位の数を下より多くした状態（反転分布）を、光や電気で無理に作ります。これは、ボルツマン分布にしたがわない、熱のつり合いからはずれた状態です（『原子』）。
\end{seikaku}

%% ---------------------------------------------------------------
\bksubsection{8}{ゆれる原子と固体の比熱}
\label{ph:熱-15:sec8}
\phboxhead{ゆれる原子の分配関数}
④で、原子のゆれのエネルギーは、粒 $\Delta\varepsilon$ の何個分かで数えられるとしました。粒を $j$ 個もつ状態のエネルギーは $j\Delta\varepsilon$（$j=0,\ 1,\ 2,\ \ldots$）です。粒の大きさ $\Delta\varepsilon$ は、ばねにつけたおもりのゆれの角振動数 $\omega=\sqrt{k_{\mathrm s}/m}$（6-6①）に比例します（比例の係数は、原-2・原-6で）。分配関数は、$x=e^{-\Delta\varepsilon/(k_{\mathrm B}T)}$ とおくと等比級数になり、$0<x<1$ なので7-10⑤の式が使えます。
\[Z=\sum_{j=0}^{\infty}e^{-j\Delta\varepsilon/(k_{\mathrm B}T)}=1+x+x^2+\cdots=\frac{1}{1-x}=\frac{1}{1-e^{-\beta\Delta\varepsilon}}\]
$\ln Z=-\ln(1-e^{-\beta\Delta\varepsilon})$ を $\beta$ で微分して $-1$ をかけると、
\pheq{\langle\varepsilon\rangle=-\frac{\dd\ln Z}{\dd\beta}=\frac{\Delta\varepsilon\,e^{-\beta\Delta\varepsilon}}{1-e^{-\beta\Delta\varepsilon}}=\frac{\Delta\varepsilon}{e^{\Delta\varepsilon/(k_{\mathrm B}T)}-1}}{J ＝ J}

\phboxhead{高い温度では $k_{\mathrm B}T$：等分配の法則}
粒が温度にくらべて小さいとき（$\Delta\varepsilon\ll k_{\mathrm B}T$）は、$a=\Delta\varepsilon/(k_{\mathrm B}T)$ が小さいので、$e^a\approx1+a$（5-17⑥）を使うと、
\[\langle\varepsilon\rangle\approx\frac{\Delta\varepsilon}{(1+a)-1}=\frac{\Delta\varepsilon}{a}=k_{\mathrm B}T\]
です（例題10）。ゆれる原子は、運動のエネルギーで $\tfrac{1}{2}k_{\mathrm B}T$、ばねの位置エネルギーで $\tfrac{1}{2}k_{\mathrm B}T$ をもち、あわせて $k_{\mathrm B}T$ です。熱-14 ④ の\textbf{エネルギー等分配の法則}（エネルギーの2乗の項1つあたり $\tfrac{1}{2}k_{\mathrm B}T$）が、高い温度の極限として出てきました。

\phboxhead{アインシュタインの固体}
固体の原子 $N$ 個は、それぞれ $x$・$y$・$z$ の3つの向きにゆれるので、ゆれは $3N$ 個あります。アインシュタインは1907年に、どのゆれも同じ $\Delta\varepsilon$ をもつと考えて、固体の比熱を計算しました。ゆれ1つの $\langle\varepsilon\rangle$ を、⑦と同じように $x=\Delta\varepsilon/(k_{\mathrm B}T)$ とおいて温度で微分すると、
\[\frac{\dd\langle\varepsilon\rangle}{\dd T}=\frac{\dd}{\dd T}\,\frac{\Delta\varepsilon}{e^x-1}=\frac{\Delta\varepsilon\,e^x}{(e^x-1)^2}\times\frac{\Delta\varepsilon}{k_{\mathrm B}T^2}=k_{\mathrm B}\frac{x^2e^x}{(e^x-1)^2}\]
です。$U=3N\langle\varepsilon\rangle$ なので、1 mol（$N=N_{\mathrm A}$、$N_{\mathrm A}k_{\mathrm B}=R$）の比熱は
\pheq{C_V=3R\,\frac{x^2e^x}{(e^x-1)^2},\qquad x=\frac{\theta_{\mathrm E}}{T},\quad \theta_{\mathrm E}=\frac{\Delta\varepsilon}{k_{\mathrm B}}}{J/(mol·K) ＝ J/(mol·K)}
\label{ph:熱-15:einstein}となります。$\theta_{\mathrm E}$ を\textbf{アインシュタイン温度}といいます。\bkidx{あいんしゅたいんのこたい}{アインシュタインの固体}

\begin{phflow}[ゆれる原子の分配関数から]
\phflowrow{$\langle\varepsilon\rangle=\dfrac{\Delta\varepsilon}{e^{\Delta\varepsilon/(k_{\mathrm B}T)}-1}$}{高い温度\\$\Delta\varepsilon\ll k_{\mathrm B}T$}{$\langle\varepsilon\rangle=k_{\mathrm B}T$\\（等分配：熱-14 ④）}
\phflowrow{$C_V=3R\,\dfrac{x^2e^x}{(e^x-1)^2}$}{高い温度\\$T\gg\theta_{\mathrm E}$（$x\to0$）}{$C_V=3R\approx\qty{25}{J/(mol.K)}$\\（デュロン・プティ）}
\end{phflow}

\begin{bkfigure}
\begin{tikzpicture}
\begin{axis}[bkplot, width=112mm, height=46mm, xmin=0, xmax=3.2, ymin=0, ymax=1.15, xtick={0.23,0.5,1,2,3}, xticklabel style={/pgf/number format/fixed}, ytick={0.25,0.5,0.75,1}, yticklabel style={/pgf/number format/fixed}, xlabel={$T/\theta_{\mathrm E}$}, ylabel={$C_V/(3R)$}]
  \addplot[bkthin, dashed, domain=0:3.1] {1};
  \addplot[bkcurve, domain=0.05:3.1, samples=140] {(1/x)^2*exp(1/x)/(exp(1/x)-1)^2};
  \fill[bkRed] (axis cs:0.2292,0.2452) circle (1.8pt);
  \draw[bkRed!80!black, line width=0.6pt] (axis cs:0.2292,0.2452) -- (axis cs:0.45,0.18);
  \node[font=\sffamily\footnotesize, text=bkRed!80!black, anchor=west] at (axis cs:0.46,0.18) {ダイヤモンド 298 K（測った値 約\qty{6.1}{J/(mol.K)}）};
  \node[font=\sffamily\footnotesize, anchor=south east] at (axis cs:3.1,1.02) {デュロン・プティの $3R$};
\end{axis}
\end{tikzpicture}
\bkfigcaption{アインシュタインの固体の比熱。高温で $3R$ に近づき、低温で0になる。ダイヤモンドの点は、測った値を $3R$ でわったもの（$\theta_{\mathrm E}=\qty{1300}{K}$ とした）。}\label{ph:熱-15:fig-einstein}
\end{bkfigure}
多くの金属は、室温で $T$ が $\theta_{\mathrm E}$ より大きく、比熱はほぼ $3R$ です（デュロン・プティの法則。熱-14 ④）。ところが、ダイヤモンドの比熱は、室温（298 K）で約\qty{6.1}{J/(mol.K)}しかなく、$3R$ の4分の1ほどです。$\theta_{\mathrm E}$ を約\qty{1300}{K}とすると、$T/\theta_{\mathrm E}=0.23$ で、式から $C_V\approx\qty{6.2}{J/(mol.K)}$ となり、よく合います。古い理論（等分配）では説明できなかった「比熱が温度で変わる」ことを、エネルギーの粒の考えで説明できたことは、エネルギーがとびとびの値をとるという考えが、光だけでなく物にもあてはまることを示しました。

\begin{phbaisuru}[4倍]{ばね定数 $k_{\mathrm s}$}
原子どうしをつなぐばね（結びつき）が強いと、$\omega=\sqrt{k_{\mathrm s}/m}$ は $\sqrt4=2$ 倍、$\Delta\varepsilon$ と $\theta_{\mathrm E}$ も2倍になります。同じ温度では $T/\theta_{\mathrm E}$ が半分になり、比熱は小さくなります。たとえば $T=\theta_{\mathrm E}$ で $0.92\times3R$ だった物は、$T=\theta_{\mathrm E}/2$ の $0.72\times3R$ になります。ダイヤモンドは、結びつきがとても強く、原子（炭素）が軽いので、$\theta_{\mathrm E}$ が大きいのです。
\end{phbaisuru}

\begin{seikaku}
本物の固体をとても低い温度まで冷やすと、比熱は $T^3$ に比例して0に近づき、アインシュタインの式（$e^{-\theta_{\mathrm E}/T}$ のように急に0になる）より、ゆっくり下がります。固体の中のゆれは、全部同じ $\omega$ ではなく、長い波から短い波まで、いろいろな $\omega$ の波でできているからです。これを取り入れたのがデバイの理論（1912年）で、原-8 で学びます。また、『原子』では、いちばん低いエネルギーが0ではなく $\tfrac{1}{2}\Delta\varepsilon$ であることも学びますが、比熱は変わりません。
\end{seikaku}

%% ---------------------------------------------------------------
\bksubsection{9}{この先へ：フェルミ分布とボース分布}
\label{ph:熱-15:sec9}
⑥では、1つの棚に何個でも分子が入れるとしました。ところが、電子は、1つの状態に1個までしか入れません（原-7）。光の粒（光子）や⑧のゆれの粒は、反対に、1つの状態にいくつでも入れて、しかも粒どうしを区別できません。このような粒子では、エネルギー $\varepsilon$ の1つの状態に入っている平均の数が、次の形になります。
\[\text{フェルミ分布}\ \ \frac{1}{e^{(\varepsilon-\mu)/(k_{\mathrm B}T)}+1},\qquad \text{ボース分布}\ \ \frac{1}{e^{(\varepsilon-\mu)/(k_{\mathrm B}T)}-1}\]
$\mu$ は\textbf{化学ポテンシャル}といい、粒子を1個加えるときに要るエネルギーです（熱-10）。⑥の $\alpha$ とは $\mu=-\alpha k_{\mathrm B}T$ の関係があります。$e^{(\varepsilon-\mu)/(k_{\mathrm B}T)}$ が1よりずっと大きい所では、どちらも $e^{-(\varepsilon-\mu)/(k_{\mathrm B}T)}$ となり、ボルツマン分布にもどります（図\ref{ph:熱-15:fig-fb}）。⑧のゆれる原子では、粒の平均の数は
\[\frac{\langle\varepsilon\rangle}{\Delta\varepsilon}=\frac{1}{e^{\Delta\varepsilon/(k_{\mathrm B}T)}-1}\]
で、これは $\mu=0$ のボース分布そのものです。
\begin{bkfigure}
\begin{tikzpicture}
\begin{axis}[bkplot, axis lines=left, width=104mm, height=46mm, xmin=-3, xmax=3.4, ymin=0, ymax=2.1, xtick={-3,-2,-1,0,1,2,3}, ytick={0,0.5,1,1.5,2}, xlabel={$(\varepsilon-\mu)/(k_{\mathrm B}T)$}, ylabel={平均の数}]
  \addplot[bkcurve, domain=-3:3.2, samples=100] {1/(exp(x)+1)};
  \addplot[bkcurve2, dashed, domain=0.38:3.2, samples=100] {1/(exp(x)-1)};
  \addplot[bkGreen!70!black, line width=1.1pt, dotted, domain=-0.74:3.2, samples=100] {exp(-x)};
  \node[font=\sffamily\footnotesize, text=bkBlue!80!black, anchor=south west] at (axis cs:-2.9,1.0) {フェルミ（実線）};
  \node[font=\sffamily\footnotesize, text=bkOrange!75!black, anchor=west] at (axis cs:0.62,1.75) {ボース（破線）};
  \node[font=\sffamily\footnotesize, text=bkGreen!50!black, anchor=south west] at (axis cs:1.5,0.5) {ボルツマン（点線）};
\end{axis}
\end{tikzpicture}
\bkfigcaption{3つの分布。$\varepsilon-\mu$ が $k_{\mathrm B}T$ の数倍より大きい所では、3つともほぼ重なる。フェルミ分布は1をこえない。}\label{ph:熱-15:fig-fb}
\end{bkfigure}
金属や半導体の中の電子はフェルミ分布、熱した物から出る光（熱ふく射。熱-13）はボース分布にしたがいます。どちらも原-2・原-8で使います。

\begin{phmada}[ガラスは、なぜ固まる？]
窓ガラスは、分子が液体のように乱れて並んだまま、固まっています。とけたガラスを冷やすと、ある温度のあたりで、流れにくさ（粘度）が何けたも急に大きくなり、固体のようにふるまいます（ガラス転移）。結晶ができるときのような、はっきりした「並び方の変わり目」は見つかっていません。なぜ固まるのかは、固体の物理の大きな未解決問題の1つとされ、1977年にノーベル物理学賞を受けたアンダーソンは、1995年に「固体の理論で最も深く、最も興味深い未解決の問題は、おそらくガラスとガラス転移の性質だろう」と書きました。今も研究が続いています。

また、この章のすべての土台である等確率の原理（①）が、なぜ成り立つのかについても、いくつかの考え方があります。分子の運動の法則から、どんな物でも成り立つことを示すことは、まだできていません。それでも、等確率の原理から計算した比熱や圧力は、実験とよく合うので、物理学者はこの原理を使い続けています。
\end{phmada}

\begin{phsaki}{量子の統計・相転移・情報のエントロピー}
⑨のフェルミ分布とボース分布を使うと、金属の電子の比熱、半導体の中の電子の数、熱した物が出す光の色（プランクの式）が計算できます（『原子』）。また、磁石を熱すると、ある温度で急に磁石でなくなるように、たくさんの粒子がいっせいに並び方を変える現象（相転移）を、分配関数から調べる分野があり、専門の本で学びます。通信の分野では、シャノン（1948年）が、情報の量を $-\sum P_i\ln P_i$ の形の「エントロピー」で表しました。等確率のとき（$P_i=1/\Omega$）、これは $\ln\Omega$ になり、ボルツマンのエントロピーと同じ形です。
\end{phsaki}

\begin{senmonsho}[統計力学の専門書の書き方]
\begin{bkhonbun}
孤立系において，エネルギー $E$ をもつ微視的状態の数を $W(E)$ とすると，エントロピーは $S=k_{\mathrm B}\ln W$ で与えられる．温度 $T$ の熱浴と接した系がエネルギー $E_i$ の状態にある確率は $P_i=e^{-\beta E_i}/Z$（$\beta=1/k_{\mathrm B}T$）であり，$Z=\sum_ie^{-\beta E_i}$ を分配関数という．ヘルムホルツの自由エネルギーは $F=-k_{\mathrm B}T\ln Z$ で与えられる．
\end{bkhonbun}
\bkyaku{「孤立系」は、まわりとエネルギーをやりとりしない物。「微視的状態」はミクロな状態で、$W$ はこの本の $\Omega$（この本では $W$ は仕事なので $\Omega$ と書いた）。「熱浴」は、⑥のもっと知りたいの「温度の決まった大きなまわり」。$E_i$ はこの本の $\varepsilon_i$。この確率の形を\textbf{カノニカル分布}ともいう。$F$ は熱-10 の自由エネルギー $U-TS$ で、$Z$ から $F$ がわかると、圧力・エントロピーなども計算できる。}
\end{senmonsho}

\begin{taisetsu}
\def\tsl#1{\makebox[9.4\zw][l]{\sffamily\bfseries\small\color{\bkTT}#1}}
\tsl{状態の数}$\Omega$：1つのマクロな状態に入るミクロな状態の数。$P\propto\Omega$\\[3pt]
\tsl{ボルツマンの関係}$S=k_{\mathrm B}\ln\Omega$。温度は $\dfrac{1}{T}=\dfrac{\Delta S}{\Delta U}$\\[3pt]
\tsl{スターリング}$\ln n!\approx n\ln n-n$\\[3pt]
\tsl{ボルツマン分布}$P_i=\dfrac{e^{-\varepsilon_i/(k_{\mathrm B}T)}}{Z}$、$Z=\sum_ie^{-\varepsilon_i/(k_{\mathrm B}T)}$\\[3pt]
\tsl{平均のエネルギー}$\langle\varepsilon\rangle=-\dfrac{\dd\ln Z}{\dd\beta}$（$\beta=1/(k_{\mathrm B}T)$）\\[3pt]
\tsl{2準位・ゆれ}比 $e^{-\Delta\varepsilon/(k_{\mathrm B}T)}$。ゆれは $\langle\varepsilon\rangle=\dfrac{\Delta\varepsilon}{e^{\Delta\varepsilon/(k_{\mathrm B}T)}-1}$、高温で $k_{\mathrm B}T$
\end{taisetsu}

\begin{kotoba}
\small
\bkword{ミクロな状態}{1つ1つの分子がどこにいて、どんなエネルギーをもつかまで決めた状態。}
\bkword{マクロな状態}{左にいる数・全体のエネルギーなど、全体の量だけで見た状態。}
\bkword{状態の数 $\Omega$}{1つのマクロな状態に入る、ミクロな状態の数。}
\bkword{等確率の原理}{条件が同じなら、どのミクロな状態も同じ確からしさで起こるという仮定。}
\bkword{ボルツマンの関係}{エントロピーは状態の数の対数に比例する：$S=k_{\mathrm B}\ln\Omega$。}
\bkword{スターリングの近似}{大きな $n$ で $\ln n!\approx n\ln n-n$。}
\bkword{ボルツマン分布}{温度 $T$ で、エネルギー $\varepsilon$ の状態にいる確率は $e^{-\varepsilon/(k_{\mathrm B}T)}$ に比例する。}
\bkword{分配関数 $Z$}{ボルツマン因子の和。平均のエネルギーや比熱を出すもとになる。}
\end{kotoba}

\bkkigocards{\bkkigocard{$\Omega$}{オメガ}{状態の数。$S=k_{\mathrm B}\ln\Omega$。専門書の $W$（この本の $W$ は仕事）}{microstates}{熱-15}}
{\bkkigocard{$Z$}{ゼット}{分配関数。ボルツマン因子の和。熱-10 の圧縮係数とは別}{partition function}{熱-15}}

%% ======================== 例題 ========================
\Needspace{12\baselineskip}\bksection{reidai}{例題でたしかめよう}
\label{ph:熱-15:ex1}
\begin{phreidai}[52枚のトランプの並べ方]
52枚のトランプの並べ方は、何通りありますか。また、よく切ったとき、新品のときと同じ順（マークごとに A から K）にちょうどなる確率を求めましょう。
\phbamen
\bkhatenaref のトランプです。カードを上から1列に並べた順番を数えます。
\phmushi
52枚はどれもちがうカードで、区別できるとします。「よく切る」と、どの並べ方も同じ確からしさで出るとします（等確率）。
\phzu
上から1枚め、2枚め、…の位置に、残っているカードから1枚ずつ置いていきます。並べ方の数を $\Omega$、確率を $P$ とします。
\phriyu
順番に選ぶ選び方の数は、それぞれの段の選び方の数のかけ算です（3-17②の積の法則）。どの並べ方も同じ確からしさなので、1つの並べ方の確率は $1/\Omega$ です。
\phtateru
\[\Omega=52\times51\times50\times\cdots\times2\times1=52!,\qquad P=\frac{1}{52!}\]
\phtoku
表計算の FACT(52) や、Python の \texttt{math.factorial(52)} で計算できます。手で見当をつけるには、常用対数の和を使います（3-12）。
\[\log_{10}52!=\log_{10}1+\log_{10}2+\cdots+\log_{10}52\approx67.907\]
\[\Omega=10^{67.907}\approx\num{8.07e67}\ \text{通り},\qquad P=\frac{1}{\num{8.07e67}}\approx\num{1.2e-68}\]
\phtashikame
スターリングの近似（⑤のせいかくに言うとの形）で計算すると、
\[\ln52!\approx52\ln52-52+\frac{1}{2}\ln(2\pi\times52)\approx156.36\]
常用対数にすると $156.36\div2.303\approx67.91$ で、同じ値です。\\
大きさ：3-17のもっと知りたいの「約 $8\times10^{67}$ 通り」と合います。
\kotae 約 $\num{8.07e67}$ 通り。確率は約 $\num{1.2e-68}$
\end{phreidai}
\begin{jibun}
1つのマークの13枚（A〜K）の並べ方は何通りですか。また、よく切ったとき、A、2、3、…、K の順にちょうどなる確率を求めましょう。
\bkjans{$13!=\num{6227020800}\approx\num{6.2e9}$ 通り。確率は約 $\num{1.6e-10}$}
\end{jibun}

\label{ph:熱-15:ex2}
\begin{phreidai}[コイン100枚で、表がちょうど50枚]
コインを100枚投げます。表がちょうど50枚になる並び方は何通りですか。全部表になる並び方とくらべましょう。また、表がちょうど50枚になる確率を求めましょう。
\phbamen
①の分子が左右どちらにいるかと同じ数え方です。
\phmushi
コインは表と裏が同じ確からしさで出て、1枚1枚は区別できるとします。
\phzu
コインの数 $N=100$、表の数 $k=50$。並び方の数 $\Omega(50)$、全部の並び方 $2^{100}$。
\phriyu
「100枚のうち、どの50枚が表か」の選び方なので組合せ ${}_{100}\mathrm{C}_{50}$ です（①）。確率は $\Omega/2^N$。
\phtateru
\[\Omega(50)={}_{100}\mathrm{C}_{50}=\frac{100!}{50!\times50!},\qquad P=\frac{\Omega(50)}{2^{100}}\]
\phtoku
Python の \texttt{math.comb(100, 50)} などで計算すると、
\[\Omega(50)\approx\num{1.01e29},\qquad 2^{100}\approx\num{1.27e30},\qquad P\approx\frac{\num{1.01e29}}{\num{1.27e30}}\approx0.080\]
全部表は1通りなので、表50枚の並び方は、全部表の約 $10^{29}$ 倍です。
\phtashikame
⑤の式で見積もると、$\ln{}_{100}\mathrm{C}_{50}\approx100\ln2=69.3$ で、$e^{69.3}\approx10^{30}$。けたは近く、少し大きめに出ます（⑤の近似で $\tfrac{1}{2}\ln(2\pi n)$ をすてた分）。\\
極端な場合：ちょうど50枚の確率は0.08で、それほど大きくありません。けれども、表が40〜60枚のどれかになる確率は0.96です（①）。
\kotae 約 $\num{1.0e29}$ 通り（全部表の約 $10^{29}$ 倍）。確率は約0.080
\end{phreidai}
\begin{jibun}
コインを20枚投げます。表がちょうど10枚になる並び方は何通りですか。また、その確率を求めましょう（$2^{20}=\num{1048576}$）。
\bkjans{${}_{20}\mathrm{C}_{10}=\num{184756}$ 通り。確率は約0.18}
\end{jibun}

\label{ph:熱-15:ex3}
\begin{phreidai}[分子が全部、片方に集まる確率]
箱の中に分子が100個あります。ある瞬間に、100個全部が箱の左半分にいる確率を求めましょう。また、1秒に $10^9$ 回、箱の中をのぞけるとすると、全部が左にいる瞬間を1回見るのに、平均で何年かかりますか。
\phbamen
気体が自然に片側に集まり、反対側が真空になることはあるか、という問いです。
\phmushi
分子は互いにじゃまをせず、1個1個が左右どちらにいるかは同じ確からしさで、ほかの分子と関係なく決まるとします。のぞくたびの結果も、前の結果と関係しないとします（本当は、続けてのぞくと似た並びが出ます）。
\phzu
分子の数 $N=100$。1個が左にいる確率 $\tfrac{1}{2}$。全部左の確率 $P$。のぞく回数は1秒に $10^9$ 回。
\phriyu
1個ずつ関係なく決まるので、確率はかけ算です（3-17②）。平均の回数は、確率の逆数です（3-17④の期待値の考え）。
\phtateru
\[P=\Bigl(\frac{1}{2}\Bigr)^{100}=2^{-100},\qquad \text{平均の回数}=\frac{1}{P}=2^{100}\]
\phtoku
$P=2^{-100}\approx\num{7.9e-31}$。平均の回数は $2^{100}\approx\num{1.27e30}$ 回。1秒に $10^9$ 回なので、
\[\frac{\num{1.27e30}}{10^9}\ \text{s}=\qty{1.27e21}{s},\qquad \frac{\qty{1.27e21}{s}}{\qty{3.16e7}{s/年}}\approx\num{4e13}\ \text{年}\]
\phtashikame
宇宙の年齢（約 $\num{1.4e10}$ 年）の約3000倍です。分子100個でもこうなので、$10^{23}$ 個の本物の気体が片側に集まることは、事実上ありません。\\
極端な場合：分子が1個なら確率は $\tfrac{1}{2}$ で、いつでも起こります。
\kotae 約 $\num{7.9e-31}$。平均で約 $\num{4e13}$ 年
\end{phreidai}
\begin{jibun}
分子が20個のとき、全部が左半分にいる確率を求めましょう。
\bkjans{$2^{-20}\approx\num{9.5e-7}$（約100万回に1回）}
\end{jibun}

\label{ph:熱-15:ex4}
\begin{phreidai}[自由膨張のエントロピーを数えて求める]
理想気体1 mol が、箱の左半分に入っています。仕切りを外すと、気体は右半分（真空）へ広がり、体積が2倍になりました。温度は変わりません。エントロピーの増えを、状態の数から求めましょう。
\phbamen
熱-6 ④ の自由膨張を、並び方の数で考えます。
\phmushi
分子どうしの力を無視します（理想気体）。温度が変わらないので、分子の速さのちらばりは変わらず、変わるのは分子のいられる場所だけとします。
\phzu
分子の数 $N=N_{\mathrm A}=\num{6.02e23}$。広がる前と後の状態の数 $\Omega_1$、$\Omega_2$。
\phriyu
1個の分子が入れる場所の数は2倍になり、$N$ 個では $2^N$ 倍です（④）。$S=k_{\mathrm B}\ln\Omega$ の差をとります。
\phtateru
\[\Delta S=k_{\mathrm B}\ln\Omega_2-k_{\mathrm B}\ln\Omega_1=k_{\mathrm B}\ln\frac{\Omega_2}{\Omega_1}=k_{\mathrm B}\ln2^{N}=Nk_{\mathrm B}\ln2\]
\phtoku
$Nk_{\mathrm B}=N_{\mathrm A}k_{\mathrm B}=R=\qty{8.314}{J/K}$（1 mol）なので、
\[\Delta S=8.314\times\ln2=8.314\times0.6931\approx\qty{5.76}{J/K}\]
\phtashikame
熱-6 ④ で、温度を保ってゆっくり広げる道の $\int\delta Q/T$ から求めた値 $R\ln2$ と同じです。\\
単位：$k_{\mathrm B}$ の J/K に、ただの数 $\ln$ をかけるので J/K。\\
2 mol なら $N$ が2倍なので、$\Delta S$ も2倍です。
\kotae 約\qty{5.76}{J/K}（$R\ln2$）
\end{phreidai}
\begin{jibun}
理想気体1 mol の体積が、温度を変えずに3倍に広がりました。エントロピーの増えを、状態の数から求めましょう。
\bkjans{$\Delta S=Nk_{\mathrm B}\ln3=R\ln3\approx\qty{9.13}{J/K}$}
\end{jibun}

\label{ph:熱-15:ex5}
\begin{phreidai}[スターリングの近似は、どれくらい当たる？]
$n=10$ と $n=100$ で、$\ln n!$ の本当の値と、スターリングの近似 $n\ln n-n$ をくらべ、ずれ（差÷本当の値）を求めましょう。
\phbamen
⑥で $\ln N!$ を近似するとき、どれくらい信じてよいかを確かめます。
\phmushi
計算の丸めのずれは無視します（4けたで計算）。
\phzu
本当の値 $\ln n!$、近似 $s_n=n\ln n-n$、ずれ $(\ln n!-s_n)/\ln n!$。
\phriyu
$\ln n!=\ln1+\cdots+\ln n$ は表計算で LN を $n$ 行たせば求まります（⑤）。近似は⑤の式です。
\phtateru
\[\ln10!=\ln\num{3628800},\qquad s_{10}=10\ln10-10,\qquad s_{100}=100\ln100-100\]
\phtoku
$\ln10!=\ln\num{3628800}\approx15.10$、$s_{10}=10\times2.3026-10\approx13.03$。ずれは $(15.10-13.03)\div15.10\approx0.14$。\\
$\ln100!\approx363.74$（LN の和）、$s_{100}=100\times4.6052-100\approx360.52$。ずれは $(363.74-360.52)\div363.74\approx0.0089$。
\phtashikame
$n$ を10倍にすると、ずれは14\,\%から0.9\,\%へ小さくなりました（⑤の表）。近似はいつも少し小さめに出ます（図\ref{ph:熱-15:fig-stirling}の長方形は、曲線の下の面積より大きい）。
\kotae $n=10$ でずれ約14\,\%、$n=100$ で約0.9\,\%
\end{phreidai}
\begin{jibun}
$n=\num{1000}$ で、$\ln n!\approx\num{5912.13}$ です。スターリングの近似とのずれを求めましょう。
\bkjans{$1000\ln1000-1000\approx\num{5907.76}$。ずれは約0.074\,\%}
\end{jibun}

\label{ph:熱-15:ex6}
\begin{phreidai}[窒素の分子は、どの高さまで多い？]
温度が15\,℃で一定の大気の中で、窒素の分子の数（1 m$^3$ あたり）が地面の $1/e$ になる高さ $H$ を求めましょう。また、富士山の山頂（\qty{3776}{m}）では、地面の何倍ですか。窒素のモル質量は\qty{28.0}{g/mol}とします。
\phbamen
空気が上空ほどうすくなるわけを、ボルツマン分布で考えます。
\phmushi
大気の温度は高さによらず一定（本当は、上空ほど低い：熱-3 ⑥）。風でまざることを無視し、窒素だけを考えます。$g$ は高さによらず\qty{9.8}{m/s^2}とします。
\phzu
高さ $z$、分子1個の質量 $m$、数 $n_V(z)$。$T=15+273.15=\qty{288.15}{K}$。
\phriyu
高さ $z$ の分子の位置エネルギーは $mgz$ で、温度が決まっているので、ボルツマン分布（⑥）が使えます。
\phtateru
\[n_V(z)=n_V(0)\,e^{-mgz/(k_{\mathrm B}T)},\qquad H=\frac{k_{\mathrm B}T}{mg}=\frac{RT}{Mg}\]
\phtoku
$M=\qty{0.0280}{kg/mol}$ を使って、
\[H=\frac{8.314\times288.15}{0.0280\times9.8}\approx\qty{8.73e3}{m}\approx\qty{8.7}{km}\]
富士山の山頂では $e^{-3776/8730}=e^{-0.433}\approx0.65$ 倍です。
\phtashikame
熱-3 ⑥ の、空気の高さの目安とほぼ同じ大きさです（空気のモル質量\qty{28.96}{g/mol}を使うと約\qty{8.4}{km}）。\\
単位：$\dfrac{\text{J/(mol·K)}\times\text{K}}{\text{kg/mol}\times\text{m/s}^2}=\dfrac{\text{J}}{\text{kg·m/s}^2}=\dfrac{\text{N·m}}{\text{N}}=$ m。\\
極端な場合：温度を上げると $H$ は大きく、重い分子ほど $H$ は小さくなります。
\kotae 約\qty{8.7}{km}。富士山の山頂で約0.65倍
\end{phreidai}
\begin{jibun}
同じ大気（15\,℃で一定）で、酸素（モル質量\qty{32.0}{g/mol}）の分子の数が $1/e$ になる高さと、富士山の山頂での割合を求めましょう。
\bkjans{$H=RT/(Mg)$ に $M=\qty{0.0320}{kg/mol}$ を入れて、約\qty{7.6}{km}。山頂で $e^{-3776/7640}$ より、約0.61倍}
\end{jibun}

\label{ph:熱-15:ex7}
\begin{phreidai}[2つの準位の数の比]
エネルギーの差が\qty{0.1}{eV}の2つの準位をもつ粒子が、300 K のまわりと熱のつり合いにあります。上の準位にいる数は、下の準位にいる数の何倍ですか。
\phbamen
原子や分子の中の、エネルギーのとびとびの状態に、温度でどれだけ上がるかを考えます。
\phmushi
ほかの準位は考えません（2準位）。粒子どうしは互いにじゃまをしないとします。
\phzu
$\Delta\varepsilon=\qty{0.1}{eV}$、$T=\qty{300}{K}$。上の数÷下の数 $=r$。
\phriyu
温度の決まったまわりとつり合っているので、ボルツマン分布（⑥）が使えます。比はボルツマン因子の比（⑦）。
\phtateru
\[r=e^{-\Delta\varepsilon/(k_{\mathrm B}T)}\]
\phtoku
まず $k_{\mathrm B}T$ を計算して、eV に直します。
\[k_{\mathrm B}T=\num{1.380649e-23}\times300=\qty{4.14e-21}{J}\]
\[\frac{\num{4.14e-21}}{\num{1.602e-19}}\approx\qty{0.0259}{eV}\]
\[\frac{\Delta\varepsilon}{k_{\mathrm B}T}=\frac{0.1}{0.0259}\approx3.87,\qquad r=e^{-3.87}\approx0.021\]
\phtashikame
単位：eV÷eV でただの数なので、$e$ の肩にのせられます（J と eV をまぜないこと）。\\
大きさ：差が $k_{\mathrm B}T$ の約4倍なので、上にいるのは50個に1個ほどです。
\kotae 約0.021倍
\end{phreidai}
\begin{jibun}
例題7の粒子を、600 K のまわりとつり合わせると、上と下の数の比は何倍になりますか。
\bkjans{$k_{\mathrm B}T\approx\qty{0.0517}{eV}$、$e^{-0.1/0.0517}\approx0.14$ 倍}
\end{jibun}

\label{ph:熱-15:ex8}
\begin{phreidai}[2準位の比熱の山はどこ？]
2準位の物の熱容量 $C_{\mathrm{th}}=Nk_{\mathrm B}\dfrac{x^2e^x}{(e^x+1)^2}$（$x=\Delta\varepsilon/(k_{\mathrm B}T)$）が、いちばん大きくなる温度と、そのときの熱容量を求めましょう。
\phbamen
⑦のショットキーの山の位置を、式から求めます。
\phmushi
2準位のほかの準位や、原子のゆれの比熱は考えません。
\phzu
$f(x)=\dfrac{x^2e^x}{(e^x+1)^2}$。山の $x$ を $x_0$ とします。温度は $k_{\mathrm B}T=\Delta\varepsilon/x$。
\phriyu
いちばん大きい所では、微分が0です（5-7）。$\ln f$ を微分すると計算が楽です（$f$ が最大なら $\ln f$ も最大）。
\phtateru
\[\ln f=2\ln x+x-2\ln(e^x+1),\qquad \frac{\dd\ln f}{\dd x}=\frac{2}{x}+1-\frac{2e^x}{e^x+1}=0\]
整理すると、$\dfrac{2}{x}=\dfrac{e^x-1}{e^x+1}$ です。
\phtoku
この式は、式の変形だけでは解けないので、表で $x$ を少しずつ変えて、両辺が等しくなる所をさがします（2-13⑤）。
\begin{bkcenter}
\begin{bktable}{colspec={Q[c,m,wd=22mm] *{4}{Q[c,m,wd=14mm]}}, rowsep=1.2pt, row{1}={bg=bkTheme!10}}
$x$ & 2.2 & 2.3 & 2.4 & 2.5 \\
$2/x$ & 0.909 & 0.870 & 0.833 & 0.800 \\
$\dfrac{e^x-1}{e^x+1}$ & 0.800 & 0.818 & 0.834 & 0.848 \\
\end{bktable}
\end{bkcenter}
両辺がほぼ等しいのは $x_0\approx2.40$ です。だから $k_{\mathrm B}T=\Delta\varepsilon/2.40\approx0.417\Delta\varepsilon$。そのときの値は
\[f=\frac{2.40^2\times e^{2.40}}{(e^{2.40}+1)^2}=\frac{5.76\times11.02}{12.02^2}\approx0.439\]
です。
\phtashikame
図\ref{ph:熱-15:fig-2jun}の山（$0.42$、$0.44$）と合います。\\
極端な場合：$x\to0$（高温）でも $x\to\infty$（低温）でも $f\to0$ なので、間に山があるはずです。
\kotae $k_{\mathrm B}T\approx0.42\Delta\varepsilon$ のとき、$C_{\mathrm{th}}\approx0.44Nk_{\mathrm B}$
\end{phreidai}
\begin{jibun}
準位の差が\qty{0.05}{eV}の2準位の物で、熱容量が山になる温度を求めましょう。
\bkjans{$k_{\mathrm B}T=0.417\times0.05=\qty{0.0209}{eV}$ より、約\qty{240}{K}}
\end{jibun}

\label{ph:熱-15:ex9}
\begin{phreidai}[アインシュタインの固体の比熱]
アインシュタインの固体の1 mol の比熱は、$T=\theta_{\mathrm E}/2$、$T=\theta_{\mathrm E}$、$T=3\theta_{\mathrm E}$ で、それぞれ $3R$ の何倍ですか。
\phbamen
固体を温めるのに要る熱が、温度でどう変わるかを調べます。
\phmushi
原子のゆれは、どれも同じ $\Delta\varepsilon$ をもつとします（アインシュタインの考え）。固体がふくらむことは無視します（体積一定）。
\phzu
$x=\theta_{\mathrm E}/T$。$T=\theta_{\mathrm E}/2$ なら $x=2$、$T=\theta_{\mathrm E}$ なら $x=1$、$T=3\theta_{\mathrm E}$ なら $x=1/3$。
\phriyu
⑧のアインシュタインの式に $x$ を入れます。
\phtateru
\[\frac{C_V}{3R}=\frac{x^2e^x}{(e^x-1)^2}\]
\phtoku
\[x=2:\quad\frac{4\times7.389}{6.389^2}\approx0.72\qquad\qquad x=1:\quad\frac{2.718}{1.718^2}\approx0.92\]
\[x=\frac{1}{3}:\quad\frac{0.1111\times1.3956}{0.3956^2}\approx0.99\]
\phtashikame
温度を上げるほど1に近づき、デュロン・プティの $3R$（熱-14 ④）になります。図\ref{ph:熱-15:fig-einstein}の曲線とも合います。\\
極端な場合：$x$ がとても大きい（低温）と、分母の $e^x$ の2乗が勝って0に近づきます。
\kotae 約0.72倍・0.92倍・0.99倍
\end{phreidai}
\begin{jibun}
アインシュタインの固体の比熱は、$T=2\theta_{\mathrm E}$ と $T=\theta_{\mathrm E}/4$ で、それぞれ $3R$ の何倍ですか。
\bkjans{$x=\tfrac{1}{2}$ で約0.98倍、$x=4$ で約0.30倍}
\end{jibun}

\label{ph:熱-15:ex10}
\begin{phreidai}[高い温度では、ゆれる原子のエネルギーは $k_{\mathrm B}T$]
ゆれる原子1つの平均のエネルギー $\langle\varepsilon\rangle=\dfrac{\Delta\varepsilon}{e^{\Delta\varepsilon/(k_{\mathrm B}T)}-1}$ が、高い温度で $k_{\mathrm B}T$ に近づくことを示しましょう。また、$\Delta\varepsilon=0.1k_{\mathrm B}T$ のとき、$\langle\varepsilon\rangle$ は $k_{\mathrm B}T$ の何倍ですか。
\phbamen
熱-14 の等分配の法則（高温で1つのゆれに $k_{\mathrm B}T$）が、⑧の式に入っているかを確かめます。
\phmushi
⑧と同じく、ゆれのエネルギーの粒は $\Delta\varepsilon$ で一定とします。
\phzu
$a=\Delta\varepsilon/(k_{\mathrm B}T)$。高い温度は $a\ll1$。
\phriyu
$a$ が小さいとき、$e^a$ は1次の近似 $e^a\approx1+a$（5-17⑥）で書けます。
\phtateru
\[\langle\varepsilon\rangle=\frac{\Delta\varepsilon}{e^a-1}=k_{\mathrm B}T\times\frac{a}{e^a-1}\]
\phtoku
$e^a-1\approx a$ なので、$\dfrac{a}{e^a-1}\approx1$、$\langle\varepsilon\rangle\approx k_{\mathrm B}T$。\\
$a=0.1$ なら、$e^{0.1}-1=0.10517$ で、$\dfrac{0.1}{0.10517}\approx0.951$ 倍です。
\phtashikame
2次の近似 $e^a\approx1+a+\tfrac{a^2}{2}$ まで使うと $\dfrac{a}{a+a^2/2}=\dfrac{1}{1+a/2}\approx1-\dfrac{a}{2}=0.95$ で、合います。\\
極端な場合：$a$ が大きい（低温）と $e^a$ が大きく、$\langle\varepsilon\rangle$ はほぼ0です。粒1個分のエネルギーももらえないからです。
\kotae $a\ll1$ で $\langle\varepsilon\rangle\approx k_{\mathrm B}T$。$\Delta\varepsilon=0.1k_{\mathrm B}T$ で約0.95倍
\end{phreidai}
\begin{jibun}
$\Delta\varepsilon=0.2k_{\mathrm B}T$ のとき、ゆれる原子1つの平均のエネルギーは $k_{\mathrm B}T$ の何倍ですか。
\bkjans{$0.2\div(e^{0.2}-1)=0.2\div0.2214\approx0.90$ 倍}
\end{jibun}

\label{ph:熱-15:ex11}
\begin{phreidai}[1ビットを消すのに要る熱]
コンピュータのメモリーの1ビットは、0か1の2つの状態のどちらかです。どちらかわからない1ビットを、必ず0にそろえる（消す）とき、300 K のまわりへ少なくともどれだけの熱を出さなければならないでしょう。
\phbamen
情報を消す操作を、状態の数で考えます（ランダウアーの原理、1961年）。
\phmushi
メモリーの中の、ビットを表す部分以外の自由さ（電子の運動など）の変化は考えず、まわりは温度 $T=\qty{300}{K}$ の大きな物とします。
\phzu
消す前のビットの状態の数 $\Omega_1=2$、消したあと $\Omega_2=1$。まわりへ出す熱 $Q$。
\phriyu
ビットのエントロピーは $k_{\mathrm B}\ln\Omega$（④）で減ります。全体のエントロピーは減らせない（熱-6）ので、その分をまわりのエントロピーの増え $Q/T$ でうめなければなりません。
\phtateru
\[\Delta S_{\text{ビット}}=k_{\mathrm B}\ln1-k_{\mathrm B}\ln2=-k_{\mathrm B}\ln2,\qquad \frac{Q}{T}+\Delta S_{\text{ビット}}\geqq0\]
\phtoku
\[Q\geqq k_{\mathrm B}T\ln2=\num{1.38e-23}\times300\times0.693\approx\qty{2.9e-21}{J}\]
\phtashikame
大きさ：とても小さな熱です。今のコンピュータは、1ビットを処理するのに、これよりずっと大きなエネルギーを使って熱を出しています。この限界に近い熱を、小さな粒を使った実験で測った報告があります（2012年）。\\
単位：J/K×K＝J。
\kotae 少なくとも約\qty{2.9e-21}{J}
\end{phreidai}
\begin{jibun}
1 GB（$8\times10^9$ ビット）のメモリーを、300 K で全部消すのに要る熱の、いちばん小さい値を求めましょう。
\bkjans{1ビットあたり $\qty{2.87e-21}{J}$ の $8\times10^9$ 倍で、約 $\qty{2.3e-11}{J}$}
\end{jibun}

\label{ph:熱-15:ex12}
\begin{phreidai}[リフルシャッフルは、少なくとも何回？]
トランプを2つの山に分け、両手の親指で少しずつ落として1つにまぜる切り方を、リフルシャッフルといいます。52枚のどの並べ方にもなりうるには、少なくとも何回切る必要がありますか。
\phbamen
\bkhatenaref の「よく切る」には、何回切ればよいかを考えます。
\phmushi
切り方のくせ（何枚ずつ落ちるか）は考えず、1回で行ける並べ方の数の上限だけを数えます。
\phzu
\noindent\begin{minipage}[c]{0.5\linewidth}\raggedright
1回切ったあとの山の上から1枚ずつ見ると、どのカードも「左の山から来たか、右の山から来たか」の2通りです。左の山の中、右の山の中での順番は変わりません。
\end{minipage}\hfill
\begin{minipage}[c]{0.46\linewidth}\centering
\begin{tikzpicture}[x=1mm, y=1mm]
  \foreach \i/\t in {0/1,1/2,2/3} {\qxpcard{(4+\i*7,26)}{\t}{bkBlue!80!black}}
  \foreach \i/\t in {0/4,1/5,2/6} {\qxpcard{(32+\i*7,26)}{\t}{bkRed!80!black}}
  \foreach \i/\t/\c in {0/4/bkRed!80!black,1/1/bkBlue!80!black,2/2/bkBlue!80!black,3/5/bkRed!80!black,4/3/bkBlue!80!black,5/6/bkRed!80!black} {\qxpcard{(4+\i*7,4)}{\t}{\c}}
  \draw[qxp arr, bkInk] (11,20) -- (16,10.5);
  \draw[qxp arr, bkInk] (39,20) -- (34,10.5);
  \node[qxp lab] at (11,33.5) {左の山};
  \node[qxp lab] at (39,33.5) {右の山};
\end{tikzpicture}
\end{minipage}
\par\vspace{2pt}
1回の切り方の種類は、多くても $2^{52}$ 通りです。$k$ 回切ると、多くても $(2^{52})^k=2^{52k}$ 通りの並べ方にしか行けません。
\phriyu
全部の並べ方（$52!$ 通り）に行けるには、行ける並べ方の数が $52!$ 以上でなければなりません。対数をとって比べます（3-12）。
\phtateru
\[2^{52k}\geqq52!\qquad\text{だから}\qquad 52k\geqq\log_252!\]
\phtoku
$\log_252!=\log_{10}52!\div\log_{10}2=67.907\div0.30103\approx225.6$。
\[k\geqq\frac{225.6}{52}\approx4.34\qquad\text{だから}\qquad k\geqq5\ \text{回}\]
\phtashikame
これは「少なくとも」の数です。どの並べ方も同じくらいの確率になるには、もっと切る必要があります。バイヤーとダイアコニスは1992年に、52枚では約7回のリフルシャッフルで、ほぼ全部の並べ方が同じくらいの確率になることを示しました。5回は、それより少なく、矛盾しません。
\kotae 少なくとも5回（実際によくまざるのは約7回）
\end{phreidai}
\begin{jibun}
10枚のカードでは、リフルシャッフルを少なくとも何回すれば、どの並べ方にもなりうるでしょう（$\log_{10}10!\approx6.560$）。
\bkjans{$\log_210!=6.560\div0.30103\approx21.8$、$21.8\div10=2.18$ より、少なくとも3回}
\end{jibun}

\label{ph:熱-15:ex13}
\begin{phreidai}[2つの固体が粒を分け合う]
原子2個の物 A と原子2個の物 B が、エネルギーの粒4個を分け合います。A の粒の数ごとの状態の数を求め、いちばん起こりやすい分け方と、その確率を求めましょう。また、はじめに粒が全部 B にあったとき、粒はどちら向きに移りますか。
\phbamen
④の、熱い物と冷たい物を接触させる場面を、小さな数で数えます。
\phmushi
粒の大きさは A と B で同じとし、粒は A と B の間で自由に移れるとします。どの分け方（ミクロな状態）も同じ確からしさとします。
\phzu
A の粒の数 $j_{\mathrm A}$、B の粒の数 $4-j_{\mathrm A}$。原子 $N$ 個に $j$ 粒を分ける数 $\Omega(N,\ j)={}_{j+N-1}\mathrm{C}_j$。$N=2$ なら ${}_{j+1}\mathrm{C}_j=j+1$。
\phriyu
A と B はそれぞれ独立に状態をとるので、全体の状態の数は積 $\Omega_{\mathrm A}\times\Omega_{\mathrm B}$（④）。確率はそれを全部の和でわります。
\phtateru
\[\Omega_{\mathrm A}\Omega_{\mathrm B}=(j_{\mathrm A}+1)\times(4-j_{\mathrm A}+1)\]
\phtoku
\begin{bkcenter}
\begin{bktable}{colspec={Q[l,m,wd=34mm] *{5}{Q[c,m,wd=10mm]}}, rowsep=1.2pt, row{1}={bg=bkTheme!10}}
$j_{\mathrm A}$ & 0 & 1 & 2 & 3 & 4 \\
$\Omega_{\mathrm A}\times\Omega_{\mathrm B}$ & $1\times5$ & $2\times4$ & $3\times3$ & $4\times2$ & $5\times1$ \\
状態の数 & 5 & 8 & 9 & 8 & 5 \\
\end{bktable}
\end{bkcenter}
全部で $5+8+9+8+5=35$ 通り。いちばん多いのは $j_{\mathrm A}=2$（2個ずつ）で、確率は $9\div35\approx0.26$ です。はじめが $j_{\mathrm A}=0$（5通り）なら、状態の数の多いほうへ、つまり粒は B から A へ移ります。
\phtashikame
全部の数は、粒4個と仕切り3本の並べ方 ${}_7\mathrm{C}_4=35$ と一致します。\\
大きさ：原子が4個しかないので、いちばん多い分け方でも確率は0.26で、ほかの分け方もよく起こります。④のプログラムのように原子を増やすと、2個ずつ（原子1個あたり同じ）の分け方が、ほとんど確実になります。
\kotae $j_{\mathrm A}=0$ から4 で 5・8・9・8・5 通り。2個ずつが確率約0.26でいちばん起こりやすい。粒は B から A へ移る
\end{phreidai}
\begin{jibun}
原子2個の物 A と原子3個の物 B が、粒3個を分け合います。A の粒の数 $j_{\mathrm A}=0$〜3 ごとの状態の数と、いちばん起こりやすい $j_{\mathrm A}$、その確率を求めましょう（原子3個に $j$ 粒を分ける数は ${}_{j+2}\mathrm{C}_j$）。
\bkjans{$j_{\mathrm A}=0$〜3 で $1\times10$、$2\times6$、$3\times3$、$4\times1$ より 10・12・9・4 通り（全部で35）。$j_{\mathrm A}=1$ がいちばん多く、確率は約0.34}
\end{jibun}

%% ======================== まちがえやすいところ ========================
\Needspace{11\baselineskip}\bksection{tsumazuki}{まちがえやすいところをチェック}
\begin{tsumazuki}{状態の数が2倍なら、エントロピーも2倍？}
\bkng{$\Omega$ が2倍になると、$S=k_{\mathrm B}\ln\Omega$ も2倍になる。}
\bkok{$\Omega$ が2倍なら、$S$ は $k_{\mathrm B}\ln2$ だけ増えるだけ。$S$ が2倍になるのは、$\Omega$ が2乗になるとき。}
$\ln$ は、かけ算をたし算に変えます（④）。「2倍」と「たす」を取りちがえないようにします。
\end{tsumazuki}

\begin{tsumazuki}{エントロピーが減ることは、絶対に起こらない？}
\bkng{エントロピーが減る変化は、物理の法則で禁じられているので、決して起こらない。}
\bkok{起こることもあるが、分子の数が多いと、その確率がとても小さい。}
③のつぼで、玉10個なら全部が左にもどることもありました。とても小さな物（数個の分子や、水の中の小さな粒）では、エントロピーが一時的に減るゆれも観測されます。熱力学の第二法則は、たくさんの粒子でできた物の「ほぼ確実」な法則です。
\end{tsumazuki}

\begin{tsumazuki}{$k_{\mathrm B}T$ の温度は、℃？}
\bkng{20\,℃なら $k_{\mathrm B}T=\num{1.38e-23}\times20$。}
\bkok{$T$ は必ず K で入れる：$k_{\mathrm B}\times293.15$。}
$k_{\mathrm B}T$ は、分子の運動のエネルギーの大きさで、絶対温度に比例します（熱-14）。0\,℃でも分子は動いているので、℃の数を入れると、まちがった値になります。
\end{tsumazuki}

\begin{tsumazuki}{温度を上げれば、上の準位のほうが多くなる？}
\bkng{とても高い温度にすれば、上の準位の数が下の準位より多くなる。}
\bkok{温度をどんなに上げても、上と下は半分ずつに近づくだけ。}
$e^{-\Delta\varepsilon/(k_{\mathrm B}T)}$ は1より小さいからです（⑦のせいかくに言うと）。上を多くするには、レーザーのように、つり合いからはずす仕掛けが要ります。
\end{tsumazuki}

\begin{tsumazuki}{eV と J をまぜる}
\bkng{$\Delta\varepsilon=\qty{0.1}{eV}$、$k_{\mathrm B}T=\qty{4.1e-21}{J}$ で、$e^{-0.1/(4.1\times10^{-21})}$。}
\bkok{どちらも eV（0.1 と 0.0259）か、どちらも J にそろえてからわる。}
$e$ の肩にのせる数は、単位のない数でなければなりません（例題7）。
\end{tsumazuki}

%% ======================== 工学ではこう使う ========================
\Needspace{8\baselineskip}\bksection{phgear}{工学ではこう使う}
\begin{phkougaku}[ボルツマン因子で、機械と部品を設計する]
\phboxhead{MRI：とても小さな数の差を読む}
病院の MRI（磁気共鳴画像）は、体の中の水素の原子核を、磁場の中の小さな磁石として使います。原子核の向きは2通りで、⑦の2準位の物です。\qty{1.5}{T}の磁場では、2つの準位の差は約 $\qty{2.6e-7}{eV}$ しかなく、体温（310 K）の $k_{\mathrm B}T$（約\qty{0.027}{eV}）の10万分の1ほどです。だから、上と下の数の差は、全体の約100万分の5しかありません。MRI は、この小さな差を電波で読みとります。磁場を2倍の\qty{3}{T}にすると、差の準位も数の差も2倍になり、はっきりした画像になります。

\phboxhead{部品の寿命を、高い温度で早く調べる}
部品が古くなる（劣化する）のも化学反応なので、その速さはボルツマン因子にしたがって、温度で大きく変わります。たとえば、アルミ電解コンデンサーは「使う温度が\qty{10}{K}上がるごとに、寿命がおよそ半分になる」という目安が、部品メーカーの資料に書かれています。メーカーは、高い温度で短い時間試験をして、ふだんの温度での寿命を、アレニウスの式（⑥の身近な例）で見積もります。

\phboxhead{情報を消すと熱が出る}
例題11のように、1ビットを消すには、少なくとも $k_{\mathrm B}T\ln2$ の熱が出ます。今のコンピュータの熱は、ほとんどがこの限界とは別の理由（電気を流して切りかえるときの損失）で出ていますが、小さく省エネルギーな計算のしくみを考えるときの、物理の下限になっています。

\begin{phketa}[熱の世界のエネルギーの大きさ]
\phketaline{300 K の $k_{\mathrm B}T$（1 mol では $RT$）}{\qty{4.1e-21}{J}＝約\qty{0.026}{eV}（$RT\approx\qty{2.5}{kJ}$）}
\phketaline{MRI（\qty{1.5}{T}）の水素の原子核の準位の差}{約 $\qty{2.6e-7}{eV}$}
\phketaline{化学反応の山（劣化・食べ物がいたむ）}{\qtyrange{0.5}{1}{eV}ほど}
\phketaline{分子の中の原子の結びつきのエネルギー}{数 eV（炭素と水素で約\qty{4}{eV}）}
\phboxhead{答えが桁で合っているか}
$\varepsilon/(k_{\mathrm B}T)$ が1より小さい所は分子が気軽に出入りでき、10をこえる所にはほとんどいません（$e^{-10}\approx\num{4.5e-5}$）。$e^{-100}$ のような答えが出たら、0が起こらないのと同じです。eV と J、℃ と K の取りちがえを疑いましょう。
\end{phketa}
\end{phkougaku}
\phstation{統計力学,固体物理,半導体工学,情報理論}{微視的状態数、カノニカル分布、分配関数、ボルツマン因子、アインシュタイン模型、ショットキー比熱、ランダウアーの原理}

\begin{column}{ボルツマンの墓}
オーストリアのウィーンの中央墓地に、ルートヴィッヒ・ボルツマン（1844〜1906年）の墓があります。胸像の上には、$S=k\log W$ という式が刻まれています（$\log$ は自然対数 $\ln$、$W$ はこの本の $\Omega$）。ボルツマンは、熱を分子の運動と確率で説明しようとしましたが、そのころは原子や分子が本当にあるかどうか、まだ確かめられておらず、強い反対も受けました。原子があることは、アインシュタイン（1905年）のブラウン運動の理論と、ペランの実験（1908年ごろ）で、広く認められるようになりました（熱-14）。
\end{column}

\begin{column}{マクスウェルの悪魔}
マクスウェルは、箱の仕切りに小さな戸をつけ、速い分子だけを右へ、遅い分子だけを左へ通す「悪魔」を考えました（1871年の本『熱の理論』）。悪魔が戸を開け閉めするだけなら、仕事をせずに、右を熱く、左を冷たくでき、エントロピーが減るように見えます。この謎は長く議論され、シラード（1929年）、ランダウアー（1961年）、ベネット（1982年）らの研究で、「悪魔が分子を見て覚えた情報を消すときに、例題11の熱が出て、全体のエントロピーは減らない」と考えられるようになりました。情報が物理の量でもあることを示す話です。
\end{column}

%% ======================== 章末問題 ========================
\bkendproblems
\bkproblemgroup{基本}{3}
\bkproblem コインを6枚投げます。表がちょうど3枚になる並び方は何通りですか。また、その確率を求めましょう。
\bkpans{${}_6\mathrm{C}_3=20$ 通り。確率は $20\div64\approx0.31$}
\bkproblem 箱の中の分子30個が、ある瞬間に全部左半分にいる確率を求めましょう。
\bkpans{$2^{-30}\approx\num{9.3e-10}$}
\bkproblem 理想気体2 mol が、温度を変えずに真空の側へ広がり、体積が2倍になりました。エントロピーの増えを、状態の数から求めましょう。
\bkpans{$N=2N_{\mathrm A}$ なので $\Delta S=2R\ln2\approx\qty{11.5}{J/K}$}
\bkproblem $\ln20!\approx42.34$ です。スターリングの近似 $n\ln n-n$ の値と、ずれ（差÷本当の値）を求めましょう（$\ln20\approx2.996$）。
\bkpans{$20\times2.996-20\approx39.91$。ずれは約5.7\,\%}
\bkproblem エネルギーの差が\qty{0.2}{eV}の2つの準位で、1000 K のとき、上の準位と下の準位の数の比を求めましょう。
\bkpans{$k_{\mathrm B}T\approx\qty{0.0862}{eV}$、$e^{-0.2/0.0862}\approx0.098$ 倍}
\bkproblem 温度15\,℃で一定の大気で、空気（モル質量\qty{28.96}{g/mol}）の分子の数が地面の $1/e$ になる高さを求めましょう。
\bkpans{$H=RT/(Mg)$ に $M=\qty{0.02896}{kg/mol}$ を入れて、約\qty{8.4}{km}}
\bkproblem 原子3個の物 A と原子3個の物 B が、エネルギーの粒4個を分け合います。A の粒の数 $j_{\mathrm A}=0$〜4 ごとの状態の数を求め、いちばん起こりやすい分け方と、その確率を求めましょう。
\bkpans{$1\times15$、$3\times10$、$6\times6$、$10\times3$、$15\times1$ より 15・30・36・30・15 通り（全部で126）。2個ずつがいちばん多く、確率は $36\div126\approx0.29$}
\bkproblemgroup{チャレンジ}{5}
\bkproblem 準位の差が\qty{0.05}{eV}の2準位の粒子 $N$ 個が、400 K のまわりとつり合っています。（1）上の準位にいる割合、（2）1個あたりの平均のエネルギー〔eV〕を求めましょう。
\bkpans{（1）\nobreak $k_{\mathrm B}T\approx\qty{0.0345}{eV}$、$1\div(e^{0.05/0.0345}+1)\approx0.19$。（2）\nobreak $0.05\times0.19\approx\qty{0.0095}{eV}$}
\bkproblem アインシュタイン温度が\qty{300}{K}の固体1 mol の比熱を、150 K と 600 K で求めましょう（$3R=\qty{24.9}{J/(mol.K)}$）。
\bkpans{150 K（$x=2$）：$0.724\times24.9\approx\qty{18}{J/(mol.K)}$。600 K（$x=\tfrac{1}{2}$）：$0.979\times24.9\approx\qty{24}{J/(mol.K)}$}
\bkproblem 同じ温度・圧力の窒素0.79 mol と酸素0.21 mol を、仕切りを外してまぜます（あわせた体積は変わらない）。それぞれの気体が、自分の体積から全体の体積へ広がったと考えて、エントロピーの増えを求めましょう。
\bkpans{窒素 $0.79R\ln(1/0.79)\approx\qty{1.55}{J/K}$、酸素 $0.21R\ln(1/0.21)\approx\qty{2.72}{J/K}$。あわせて約\qty{4.3}{J/K}}
\bkproblem 温度15\,℃で一定で、風でまざらない大気を考えます。地面での酸素と窒素の数の比を1とすると、高さ\qty{10}{km}では、この比は何倍になりますか（モル質量は窒素\qty{28.0}{g/mol}、酸素\qty{32.0}{g/mol}）。また、本物の大気の高さ10 km では、この比はほとんど変わりません。なぜでしょう。
\bkpans{$e^{-10/7.64}\div e^{-10/8.73}\approx0.85$ 倍。本物の大気は、風や対流でよくまざっているので、重さで分かれない}
\bkproblem 磁場\qty{3}{T}の MRI では、水素の原子核の2つの準位の差は約 $\qty{5.3e-7}{eV}$ です。体温（310 K、$k_{\mathrm B}T\approx\qty{0.0267}{eV}$）で、（1）上と下の数の比 $e^{-\Delta\varepsilon/(k_{\mathrm B}T)}$ を、$e^{-a}\approx1-a$ を使って求めましょう。（2）下の数と上の数の差は、全体の数の何倍ですか。
\bkpans{（1）\nobreak $a=5.3\times10^{-7}\div0.0267\approx\num{2.0e-5}$ より、比は約 $1-\num{2.0e-5}$。（2）\nobreak 下の数を1とすると、上は $1-a$、差は $a$、全体はほぼ2。差÷全体は $a/2\approx\num{1.0e-5}$（100万分の10。\qty{1.5}{T}の2倍）}

\begin{furikaeri}
\begin{checklist}
\item 並び方の数を数えて、変化がなぜ一方の向きにだけ進むのかを説明できる。\bkfurref{ph:熱-15:sec2}
\item $S=k_{\mathrm B}\ln\Omega$ で、並び方の数からエントロピーを計算できる。\bkfurref{ph:熱-15:sec4}
\item ボルツマン分布と分配関数から、2準位の系や固体の比熱を導ける。\bkfurref{ph:熱-15:sec7}
\end{checklist}
\end{furikaeri}
\bkfukidashi{mina}{トランプも空気も、「ばらばらな並び方のほうが、けたちがいに多い」から元にもどらないんだね。並び方を数えると、温度の意味まで見えてきた！}
